% PLOS Water manuscript - SUBMISSION COPY, generated by scripts/m11_build_latex.py from the Markdown sources.
% PLOS submission form: captions only (figures uploaded separately as FigN.tif), continuous line numbers, double spacing. The reading copy with figures is main.tex.
% Status: built; not cleared for submission; placeholders and author items remain.
% Template for PLoS
% Version 3.8 Apr 2026
%
% % % % % % % % % % % % % % % % % % % % % %
%
% -- IMPORTANT NOTE
%
% This template contains comments intended 
% to minimize problems and delays during our production 
% process. Please follow the template instructions
% whenever possible.
%
% % % % % % % % % % % % % % % % % % % % % % % 
%
% Once your paper is accepted for publication, 
% PLEASE REMOVE ALL TRACKED CHANGES in this file 
% and leave only the final text of your manuscript. 
% PLOS recommends the use of latexdiff to track changes during review, as this will help to maintain a clean tex file.
% Visit https://www.ctan.org/pkg/latexdiff?lang=en for info.
%
%
% IMPORTANT
% Below are a few tips to help format manuscript according to our specifications. For more tips to help reduce the possibility of formatting errors during conversion, please see our LaTeX guidelines at http://journals.plos.org/plosone/s/latex
%
% There are no restrictions on package use within the LaTeX files except that no packages listed in the template may be deleted.
%
% Please do not include colors or graphics in the text.
%
% The manuscript LaTeX source should be contained within a single file (do not use \input, \externaldocument, or similar commands).
%
% % % % % % % % % % % % % % % % % % % % % % %
%
% -- FIGURES AND TABLES
%
% Please include tables/figure captions directly after the paragraph where they are first cited in the text.
%
% DO NOT INCLUDE GRAPHICS IN YOUR MANUSCRIPT
% - Figures should be uploaded separately from your manuscript file. 
% - Figures generated using LaTeX should be extracted and removed from the PDF before submission. 
% - Figures containing multiple panels/subfigures must be combined into one image file before submission.
% For figure citations, please use "Fig" instead of "Figure".
% See http://journals.plos.org/plosone/s/figures for PLOS figure guidelines.
%
% Tables should be cell-based and may not contain:
% - spacing/line breaks within cells to alter layout or alignment
% - do not nest tabular environments (no tabular environments within tabular environments)
% - no graphics or colored text (cell background color/shading OK)
% See http://journals.plos.org/plosone/s/tables for table guidelines.
%
% For tables that exceed the width of the text column, use the adjustwidth environment as illustrated in the example table in text below.
%
% % % % % % % % % % % % % % % % % % % % % % % %
%
% -- EQUATIONS, MATH SYMBOLS, OTHER SPECIAL FORMATTING
%
% - For bold symbols, use the \boldsymbol command and not \mathbf. 
% - For inline equations, please be sure to include all portions of an equation in the math environment.  For example, x$^2$ is incorrect; this should be formatted as $x^2$ (or $\mathrm{x}^2$ if the romanized font is desired).
% - Do not include text that is not math in the math environment. For example, CO2 should be written as CO\textsubscript{2} instead of CO$_2$.
% - Avoid capturing simple text as equations, for example $<$.
% - For inline equations, please do not include punctuation (commas, etc) within the math environment unless this is part of the equation.
% - When adding superscript or subscripts outside of brackets/braces, please group using {}.  For example, change "[U(D,E,\gamma)]^2" to "{[U(D,E,\gamma)]}^2". 
% - Do not use \cal for caligraphic font.  Instead, use \mathcal{}.
% - Avoid splitting for fragmenting a single equation into multiple objects/environments, for example $f(x)$ = $g(x)$ + $h(x)$.
% - Use math environment for all equations, do not mimic using text.
% - Avoid using single letters in math mode where possible.
% - Always use matching parentheses and delimiters, avoid mixing math and text brackets.
% - Insert spaces around inline equations to ensure proper display after conversion.
% - Do not use forced numbers for display equation labeling or suppress numbering with \nonumber. 
%
% % % % % % % % % % % % % % % % % % % % % % % % 
%
% Please contact customercare@plos.org with any questions.
%
% % % % % % % % % % % % % % % % % % % % % % % %

\documentclass[10pt,letterpaper]{article}
\usepackage[top=0.85in,left=2.75in,footskip=0.75in]{geometry}

% amsmath and amssymb packages, useful for mathematical formulas and symbols
\usepackage{amsmath,amssymb}

% Use adjustwidth environment to exceed column width (see example table in text)
\usepackage{changepage}

% textcomp package and marvosym package for additional characters
\usepackage{textcomp,marvosym}

% cite package, to clean up citations in the main text. Do not remove.
\usepackage{cite}

% Use nameref to cite supporting information files (see Supporting Information section for more info)
\usepackage{nameref,hyperref}

% line numbers
\usepackage[right]{lineno}

% ligatures disabled
\usepackage[nopatch=eqnum]{microtype}
\DisableLigatures[f]{encoding = *, family = * }

% color can be used to apply background shading to table cells only
\usepackage[table]{xcolor}

% array package and thick rules for tables
\usepackage{array}

% create "+" rule type for thick vertical lines
\newcolumntype{+}{!{\vrule width 2pt}}

% create \thickcline for thick horizontal lines of variable length
\newlength\savedwidth
\newcommand\thickcline[1]{%
  \noalign{\global\savedwidth\arrayrulewidth\global\arrayrulewidth 2pt}%
  \cline{#1}%
  \noalign{\vskip\arrayrulewidth}%
  \noalign{\global\arrayrulewidth\savedwidth}%
}

% \thickhline command for thick horizontal lines that span the table
\newcommand\thickhline{\noalign{\global\savedwidth\arrayrulewidth\global\arrayrulewidth 2pt}%
\hline
\noalign{\global\arrayrulewidth\savedwidth}}


% Remove comment for double spacing
\usepackage{setspace}
\doublespacing

% Text layout
\raggedright
\setlength{\parindent}{0.5cm}
\textwidth 5.25in
\textheight 8.75in

% Bold the 'Figure #' in the caption and separate it from the title/caption with a period
% Captions will be left justified
\usepackage[aboveskip=1pt,labelfont=bf,labelsep=period,justification=raggedright,singlelinecheck=off]{caption}
\renewcommand{\figurename}{Fig}

% Please use the included `plos2025.bst` as your BibTeX style
\bibliographystyle{plos2025}

% Remove brackets from numbering in List of References
\makeatletter
\renewcommand{\@biblabel}[1]{\quad#1.}
\makeatother



% Header and Footer with logo
\usepackage{lastpage,fancyhdr,graphicx}
\usepackage{epstopdf}
%\pagestyle{myheadings}
\pagestyle{fancy}
\fancyhf{}
%\setlength{\headheight}{27.023pt}
%\lhead{\includegraphics[width=2.0in]{PLOS-submission.eps}}
\rfoot{\thepage/\pageref{LastPage}}
\renewcommand{\headrulewidth}{0pt}
\renewcommand{\footrule}{\hrule height 2pt \vspace{2mm}}
\fancyheadoffset[L]{2.25in}
\fancyfootoffset[L]{2.25in}
\lfoot{\today}

%% Include all user-defined macros below

% Added packages (the template allows additions; none of its packages is removed).
% Tables 1 and 4 are longer than one page; PLOS allows long tables to span pages, and a float cannot.
\usepackage{longtable}
% A display is never set before the paragraph that first cites it.
\usepackage{flafter}
% ...and every display is set before the next section or subsection begins (\FloatBarrier at each heading).
\usepackage{placeins}
% URLs in the reference list are set in the text font, as in the journal's typeset references.
\urlstyle{same}
% No line break after the colon of a URL scheme, so no line ends with 'https:' (reference QA).
\makeatletter\def\UrlNoBreaks{\do\(\do\[\do\{\do\<\do\:}\makeatother
% A reference is not split across a page, so no URL is read half on one page and half on the next (M14 R5).
\usepackage{etoolbox}
\AtBeginEnvironment{thebibliography}{\interlinepenalty=10000\relax}
% Long-table captions follow the template's caption size; the reading profile below sets it smaller.
\captionsetup[longtable]{font=normalsize}
% T1 encoding so accented names and URL underscores survive copying text from the PDF (reference review R-06).
\usepackage[T1]{fontenc}
% Float placement tuned so displays stay close to their first citation without half-empty pages.
\renewcommand{\topfraction}{0.9}
\renewcommand{\bottomfraction}{0.8}
\renewcommand{\textfraction}{0.07}
\renewcommand{\floatpagefraction}{0.8}
\setcounter{topnumber}{3}
\setcounter{bottomnumber}{2}
\setcounter{totalnumber}{4}
\makeatletter\setlength{\@fptop}{0pt}\makeatother
% Neutral PDF metadata: no author, title placeholder, no build date or TeX installation details.
\hypersetup{hidelinks,pdftitle={Draft manuscript (title to be filled later)},pdfauthor={},pdfsubject={},pdfkeywords={},pdfcreator={}}
\pdfinfoomitdate=1
\pdftrailerid{}
\pdfsuppressptexinfo=-1



%% END MACROS SECTION


\begin{document}
\vspace*{0.2in}

% Title must be 250 characters or less.
\begin{flushleft}
{\Large
\textbf{Title: TO BE FILLED LATER} % Please use "sentence case" for title and headings (capitalize only the first word in a title (or heading), the first word in a subtitle (or subheading), and any proper nouns).
}
\newline
% Insert author names, affiliations and corresponding author email (do not include titles, positions, or degrees).
\\
Authors: TO BE FILLED LATER
\\
\bigskip
Affiliations: TO BE FILLED LATER
\\
\bigskip

Corresponding author email: TO BE FILLED LATER

\end{flushleft}
% Please keep the abstract below 300 words
\section*{Abstract}
% Justified in a local group; the template sets the rest of the document ragged right.
{\rightskip=0pt\parfillskip=0pt plus 1fil Drinking-water monitoring programmes must decide what to measure, where and with which instrument, often from evidence assembled for other purposes. We developed an evidence-to-action framework that tests whether evidence supports its intended use and returns, for each monitoring context (country \texttimes{} source archive \texttimes{} parameter), an action with the evidence required to reassess it. It was applied, using archived records and published evidence, to 26 contexts: 8 in India and 18 in local, not national, archives from Kwale County (Kenya) and Malawi, and its rules also to published arsenic-kit evidence (Bangladesh) and testing costs (Ghana). The rules returned eight distinct actions; whether these improve monitoring has not been tested prospectively. Separately, in India, 17 of 160 sufficiently sampled state-level parameter cells had intervals straddling an assumed 0.10 action level (fraction of records above threshold), so further records could change at least that many classifications. Across seven tasks in a multinational fitting archive, mean area under the receiver operating characteristic curve was 0.9396 under random record splits but 0.7696 under held-out spatial blocks (one split per design): permitted use follows validation design. Separately fitted models supported rank ordering in 5 of 13 external archive \texttimes{} parameter pairs, four in Indian groundwater, where ordering may sequence confirmatory sampling within those archives; calibration slopes were below 1 on all 13, and no transferred score was used as a probability without local recalibration. A documentary audit of 23 products found 891 of 1,242 product--criterion cells unreported in any located document; no product documented all 11 criteria a comparison requires, so none was named. Each action is recorded with the rule and evidence it rests on, so the monitoring use that the rules permit for each kind of evidence can be checked.\par}

\clearpage
\newgeometry{top=0.85in,left=1in,right=1in,footskip=0.75in}
% The template offsets the header and footer 2.25in into its 2.75in left margin; with the 1in margin set above that
% offset would put the footer rule and date off the page, so it is reset here.
\fancyhfoffset[L]{0pt}
\linenumbers

\FloatBarrier

\section*{Introduction}

Evidence on drinking-water quality still covers fewer people than monitoring is meant to protect. Over the Sustainable Development Goal period, the share of the global population covered by total estimates of safely managed drinking water rose from 35\% to 72\%, and such estimates became available for India for the first time only in 2025~\cite{JMP2025}. Surveillance is also less frequent than required: in urban areas, 21 of 102 countries reported drinking-water surveillance at 95--100\% of the required frequency~\cite{GLAAS2025}. Infrequent sampling can overestimate how safe water is~\cite{Bain2014}. A programme with a fixed testing budget therefore has to decide which parameters to measure, at which water points and with which instrument, on evidence that is thin and unevenly spread.

Public archives record where and when testing happened, not the risk across a population. In India, the Central Ground Water Board's current monitoring design samples every quality station once every five years before the monsoon and at least a quarter of the stations before and after the monsoon every year~\cite{CGWB-SOP-GWQDataAnalysis}. An exceedance rate computed from records collected under any such design reflects the sampling as much as the water. Measuring where the answer is already settled spends the budget without changing an action; not measuring where the evidence is ambiguous leaves an action resting on that ambiguity.

One response is to model where contamination is likely. Predictive modelling can relate the probability that a contaminant such as groundwater arsenic exceeds a guideline value to soil, climate and terrain predictors, producing maps meant to guide further testing rather than replace it~\cite{Podgorski2020}. National models for India have mapped groundwater fluoride from 12,600 measurements~\cite{Podgorski2018} and groundwater arsenic, with the hazard maps proposed to inform the prioritisation of testing~\cite{Podgorski2020-IJERPH-India}.

How such models are evaluated decides what they can be used for, and practice differs. Environmental records are spatially autocorrelated, so a split that separates rows at random estimates something other than performance in an unseen region~\cite{Meyer2019}, and revisited sites compound the problem. Ignoring dependence between nearby samples biases model assessment towards optimism~\cite{MeyerPebesma2022}; the national Indian models were evaluated on random training and test splits~\cite{Podgorski2018,Podgorski2020-IJERPH-India}. Discrimination and calibration are likewise different properties: a model can order locations well while its stated probabilities are distorted~\cite{NiculescuMizil2005}. Estimated risks can be unreliable even when discrimination is good~\cite{VanCalster2019}, so a model whose ordering holds in a new setting may still misstate probabilities there.

Technology evidence raises the same question at the next step. Technology assessment characterises what instruments can measure and what deploying them demands, from power supply in remote settings to cleaning, signal drift and cost~\cite{Andres2018}. A programme needs performance near the threshold at which it must act. A 2003 verification of one arsenic field kit, in which one technical and one non-technical operator analysed identical samples, found false-negative rates of 62\% and 33\% respectively when they read its colour charts, relative to 10~\textmu{}g/L~\cite{EPA-ETV-QuickLowRangeII-2003}. A microbial water-quality test has been costed at 21.0 \textpm{} 11.3 USD on average across 18 African monitoring institutions in six countries, once consumables, equipment, labour and logistics are counted, with the \textpm{} not defined in the source~\cite{Delaire2017}.

Decision analysis makes trade-offs between objectives explicit and can carry uncertainty through to the comparison of alternatives~\cite{Schuwirth2018}; probabilistic forms report how probable each rank order is~\cite{Broekhuizen2015}. A ranking method returns a leader whether or not the criteria behind it are documented. Machine-learning models always make a prediction unless they are designed to abstain~\cite{Hendrickx2024}. Neither, by default, reports that the evidence is too thin for the intended use.

Monitoring decisions therefore depend on several bodies of evidence, each of which can be strong enough for one use but not the next. Part of this difficulty is recognised: a review of machine-learning prediction of United States drinking-water contaminants named accounting for spatial autocorrelation in model training among the remaining methodological challenges~\cite{Hu2023}. We developed an evidence-to-action framework that connects record harmonisation, validation design, external transfer and documentary appraisal of instruments. At each step it tests whether the evidence supports the use a programme would make of it, and for each monitoring context it returns an action together with the evidence required to reassess that action.

We applied the framework to archived records and published evidence: India as the principal application, local archives from Kwale County in Kenya and from Malawi, and decisions in Bangladesh and Ghana that rest on published studies alone. The applications address four questions a programme faces: where further records could change a monitoring classification; how much validation design changes the performance estimate it would rely on, and whether a model's ordering and probabilities carry over to groundwater archives that took no part in fitting; what records alone, or published evidence alone, support where no model is used; and whether the documentation of named instruments is sufficient to compare them. Across these applications, the rules returned an action and the evidence required to reassess it for every context, by design rather than as a measure of accuracy, separated rank use from probability use where models transferred, and named no product because the located documentation did not support a comparison; whether the actions improve monitoring was not tested.

\FloatBarrier

\section*{Materials and methods}

\subsection*{Evidence-to-action framework}

The framework takes three bodies of evidence (an accepted monitoring archive, the output of any predictive model, and the documentary evidence about candidate instruments) and returns two things for each monitoring context, a country \texttimes{} source archive \texttimes{} parameter combination (Fig~\ref{fig:workflow}). The first is an action, which names what a programme should do next rather than scoring a place. Actions run from acquiring a baseline or repairing unusable records, through prioritising confirmatory sampling, to direct or targeted reference measurement and sentinel monitoring paired with chemistry; \nameref{S2_Table} defines all eight. The second return is the evidence required to reassess the action: the additional evidence that could move the context to a different action, stated as a kind of evidence rather than as a sample size, cost or optimal design. Three gates stand between the inputs and the output: the validation gate at step 6 of the eight-step workflow, the permission gate at step 7 and the evidence gate at step 8.\label{callout:fig:workflow}

\begin{figure}[!ht]
\caption{\textbf{Evidence-to-action framework.} Eight steps connect separate evidence strands to monitoring actions through three gates: validation (step 6), permission to use ordering or probability (step 7) and product documentation (step 8). Each output is a kind of action, stated with the evidence required to reassess it. MCDA, multi-criteria decision analysis; CGWB, Central Ground Water Board.}
\label{fig:workflow}
\end{figure}

The validation gate settles in which role a model may enter a decision: ordering locations so that confirmatory sampling is sequenced, supplying a probability to compare with an action level, or neither. The role follows from the validation design, not from the size of a reported metric. Monitoring records are not exchangeable rows, because sites are revisited and nearby observations share geology, land cover and often a sampling programme, so a random split mostly measures interpolation within the existing sampling structure. Keeping every record from a site together removes the repeat-visit component; holding out contiguous spatial blocks additionally requires the model to generalise across geography~\cite{Meyer2019}. The gate prevents a random-row figure from being read as performance in a district that contributed no training station.

Discrimination asks whether the model orders cases correctly; calibration asks whether a stated probability matches the frequency actually observed at that probability~\cite{NiculescuMizil2005}. Calibration here means model calibration, not the calibration of an instrument. A ranking tells a programme where to look first; a probability compared with a threshold tells it when to act. The permission gate therefore distinguishes three states: ordering permitted, a probability permitted, or direct measurement required. Where ordering survives a move to a new setting but calibration there has not been shown to be adequate, the model may set the order of confirmatory measurement. Where calibration on local labels (direct local measurements used to check a model's scores) has also been assessed and found adequate, or restored by local recalibration, a probability may be compared with a fixed action level. Where ordering does not survive, direct measurement is required. The gate is applied to each archive \texttimes{} parameter pair, because two archives holding the same analyte can return opposite verdicts.

Once a decision requires direct measurement it becomes a question about an instrument, governed by the evidence gate through two rules. The documentation threshold asks whether a criterion is documented for enough products for a comparison on it to be informative. The sufficiency rule asks whether a given product documents every active criterion, the 11 criteria a comparison requires (\nameref{S1_Text}), and passes its must-pass analytical, matrix, deployment and identity requirements. A product that fails the sufficiency rule is not scored: scoring an undocumented field at a middling value would rank a device that reports a poor maintenance interval below one that reports nothing. Missing documentation is unknown evidence about a product, not evidence that it performs poorly. Where no product passes the sufficiency rule, the framework abstains at product level and states the documents that would lift the abstention. Multi-criteria comparison of technology archetypes gives preference context for a class of instrument under stated weights without establishing anything about a product within that class.

Abstention is accordingly a designed output with the same role as the reject option of a selective classifier~\cite{Geifman2017}, but triggered by missing evidence rather than by expected error (see Discussion), and not a failure. Record acceptance comes first, and the model gates follow in order: a context whose records cannot be accepted returns a baseline or repair action before any model is considered, and a context whose model output fails the validation or permission gate returns direct measurement. The evidence gate is applied once, to the product catalogue, and its determination holds for every context before any instrument is named. The rules are the same whatever a product documents, so no product's values can alter the gates it is judged by.

The framework separates three monitoring functions, because devices compared before the function is fixed were never doing the same job: distributed spot screening of existing water points; a higher-frequency sentinel at selected sites, which buys temporal resolution at the cost of power, logging, drift, fouling and downtime; and confirmatory or reference measurement of a specific contaminant. The functions combine within one monitoring architecture, and each setting's action names the function it relies on.

\subsection*{Study design and unit of analysis}

The framework was applied retrospectively, to archived monitoring records and published evidence; no returned action was deployed or tested. The unit of analysis for decisions is the monitoring context, one country \texttimes{} source archive \texttimes{} parameter combination, and 26 were analysed, 8 in India, 8 in Kenya and 10 in Malawi (generation rules in \nameref{S1_Text}). The Kenyan and Malawian archives are bounded local archives, not national samples, and no result from them is a national estimate for either country. Bangladesh and Ghana are applications of the decision logic to published evidence, with no model fitted or transferred; their studies were chosen for the decision each setting poses rather than found by a systematic search, so other studies could change the evidence those applications rest on. Five populations are kept separate and never merged: 23 exact commercial products, a 33-entry technology landscape (Fig~\ref{fig:workflow}), reference laboratory methods, emerging research classes, and ten technology archetypes, classes of instrument compared by multi-criteria decision analysis (MCDA). Missing evidence was retained, not imputed.

\subsection*{Data sources and evidence roles}

Six evidence strands enter the analysis (Fig~\ref{fig:evidence-map}): the fitting archive, the Indian archives, the United States archive, the Kenyan and Malawian archives, published studies from Bangladesh and Ghana, and documentation for named products. The fitting archive is the GEMStat global water-quality database of the UNEP GEMS/Water Programme~\cite{GEMStat}: a multinational archive, predominantly of surface water, used only for model fitting and internal evaluation. Seven classification tasks were built from it, for arsenic, nitrate, conductivity, suspended solids, turbidity, fluoride and pH (thresholds in \nameref{S1_Text}).\label{callout:fig:evidence-map}

\begin{figure}[!ht]
\caption{\textbf{Evidence roles by country.} Shading shows the role of each country's evidence, not the extent of sampling: external evaluation archives in India and a four-state United States well-site extract, bounded local archives in Kenya and Malawi, and published evidence only for Bangladesh and Ghana. Boundaries: Natural Earth. Data: \nameref{S9_Data}.}
\label{fig:evidence-map}
\end{figure}

Two external arms, neither of which took part in fitting, were used to evaluate transfer. The primary arm is Indian groundwater from the Central Ground Water Board (CGWB), in two products that share a producer: National Water Data Portal records for pH, conductivity, total dissolved solids (TDS) and arsenic (the CGWB portal), and fluoride and nitrate tables recovered from the CGWB Ground Water Quality Data documents (the CGWB Year Book)~\cite{CGWB-NWDP,CGWB-GWQ2019,CGWB-GWQ2021,CGWB-GWQ2022}. CGWB samples groundwater from observation and monitoring wells, with the stated aim of assessing shallow aquifers against the Indian drinking-water standard~\cite{CGWB-AGWQR2025}, so a record describes groundwater at a monitoring well, not water as supplied. The second arm, WQP, is a 2023 well-site extract for four states from the United States Water Quality Portal~\cite{WQP} (\nameref{S1_Text}). Both external arms are groundwater, so transfer from the fitting archive crosses water-body type as well as geography. Independence rests on archive lineage and a 0.1\textdegree{} cell exclusion between training and external locations, not on a record-level audit, so the external archives are not claimed to be fully independent of the fitting archive.

The bounded local archives are three openwashdata packages published under CC BY: samples from two 2016 campaigns in Kwale County, Kenya~\cite{OWD-grdwtrsmpkwale}; a laboratory snapshot of 32 Malawian water points, 30 of them boreholes or tubewells, sampled in three rounds with different parameter suites~\cite{OWD-boreholelabdata}; and a Malawian conductivity series from loggers in monitoring wells~\cite{OWD-mwgroundwaterdata}. Surface-water samples were excluded from the Kwale archive, which records no water-point type. Acceptance checked units, identity and plausibility; no charge-balance or conductivity--TDS consistency check was applied, and no records from other Indian or Kenyan custodians, microbial records included, were used. In the Indian archive, arsenic values labelled mg/L but implausible as such, from four north-eastern states in 2017, were divided by 1,000 (\nameref{S1_Text}). Harmonisation rules, versions and accepted record counts are in \nameref{S1_Text}.

\subsection*{Thresholds and chemical basis}

An exceedance rate is a record fraction, the proportion of accepted records above a threshold, reported with its numerator and denominator. Fluoride above 1.5~mg/L and arsenic above 0.010~mg/L follow the WHO guideline values, the arsenic value being provisional~\cite{WHO2022GDWQ}. These WHO values are applied to every archive, whereas CGWB judges Indian drinking suitability against the Indian standard IS 10500~\cite{CGWB-AGWQR2025}, whose limits \nameref{S1_Text} compares with ours. Nitrate was classified above 50~mg/L as NO\textsubscript{3}, the WHO guideline value expressed as nitrate ion~\cite{WHO2022GDWQ}. The transfer tasks use nitrate above 10~mg/L as N, the value of the United States maximum contaminant level~\cite{USEPA-NPDWR}. Because 10~mg/L as N is 44.268~mg/L as NO\textsubscript{3}, and not the 50~mg/L as NO\textsubscript{3} used above, no nitrate figure is carried between the two analyses. WHO proposes no health-based guideline value for pH or TDS: 6.5--8.5 is the usual optimum pH range for supply, and water becomes increasingly unpalatable at TDS above about 1,000~mg/L~\cite{WHO2022GDWQ}; both are used here as operational ranges. Conductivity above 1,500~\textmu{}S/cm is an unverified operational project convention, not a health guideline or an Indian drinking-water limit; under the factor CGWB recommends, TDS of about 0.65 \texttimes{} conductivity~\cite{CGWB-SOP-GWQDataAnalysis}, it corresponds to about 975~mg/L TDS. The 1~NTU turbidity rule is WHO's value for supporting effective disinfection~\cite{WHO2022GDWQ}, applied here to untreated water as a project rule. Microbial results use a 100~mL basis, and faecal coliform is never relabelled as \textit{E.~coli}.

Binomial uncertainty on record-level exceedance proportions is described by Wilson 95\% intervals~\cite{Brown2001}, which do not remove spatial clustering or sampling-design bias, so rates are reported as monitoring-record signals, not as population prevalence or exposure.

\subsection*{State-level cells and local archive summaries}

The accepted CGWB records were grouped into state or Union Territory \texttimes{} parameter cells. A cell was described when it held at least 30 records at 10 or more sites. For each described cell, \nameref{S3_Table} records the median revisit interval, the median gap between distinct sample dates at one site, where a site is identified by state, district, station or well identifier and coordinates rounded to five decimal places, and the cell's support for leave-one-state-out and temporal analysis. A cell is decision-relevant when the Wilson 95\% interval on its exceedance rate straddles an action level. We set that level at 0.10 as an assumption for this analysis, not a regulatory value, and also report counts at 0.05 and 0.25. This classification allocates confirmatory effort between states and parameters and sets no action; CGWB's own response to an exceedance is at station level, monitoring affected areas as trend wells and new hotspots on a 2 km grid~\cite{CGWB-SOP-GWQDataAnalysis}. Record fractions for Kwale and the Malawi laboratory snapshot use the same thresholds. At the 67 Kwale sites sampled in both campaigns (66 for nitrate and fluoride), threshold states were compared between campaigns using site medians; for the Malawi logger series, changes of state were counted between successive daily site medians.

\subsection*{Models and validation designs}

Two information channels were modelled. Model A uses environmental and terrain covariates and the site coordinate, with nearest-station neighbour features and month added in the transfer and fold-comparison fits (\nameref{S1_Text}); Model B uses earlier measurements at the same site and needs at least two of them. Both use the histogram-based gradient-boosting classifier in scikit-learn~\cite{Pedregosa2011}, with configurations fixed in advance and no hyperparameter search. Model A's configuration is in \nameref{S1_Text}; Model B's was not retained.

Three validation designs were compared on the seven fitting tasks (Fig~\ref{fig:validation-designs}). Each split the fitting archive once into training, validation and test partitions keyed on the record (random-row), the station (station-grouped) or a contiguous 3\textdegree{} spatial block; Model A was fitted on 24 environmental, terrain and coordinate features, calibrated isotonically on the validation partition and scored once on the test partition (\nameref{S1_Text}). Block splitting keeps nearby observations together and is the design used here to address generalisation to held-out blocks within the sampled geography~\cite{Meyer2019}; no autocorrelation range was estimated to set the block size. ROC-AUC (area under the receiver operating characteristic curve), precision--recall AUC, sensitivity, specificity, calibration slope and Brier skill were interpreted separately; precision--recall AUC is read against prevalence as its no-skill baseline~\cite{Saito2015}, and a 0.5 cut-point was used diagnostically rather than as an operational default. Prevalence here is a task's base rate, the share of its records above threshold, not population prevalence. A calibration slope between 0 and 1 means predicted probabilities are too extreme, and a slope at or below 0 means they carry no usable probability ordering. Model A and Model B were compared in a third set of fits, not under the three designs: forward-chained yearly folds, each year predicted from earlier years, on records with at least two earlier same-site measurements, with DeLong tests on the pooled predictions (\nameref{S1_Text}). No local recalibration was performed, because it needs held-out predictions that were not retained.\label{callout:fig:validation-designs}

\begin{figure}[!ht]
\caption{\textbf{Three validation designs.} How each design splits monitoring records into training and held-out sets: by random rows, by whole stations or by contiguous 3\textdegree{} spatial blocks, with one partition per design. Positions are schematic and carry no geographic meaning.}
\label{fig:validation-designs}
\end{figure}

A mechanistic ablation tested whether aquifer, soil, land-cover and population variables can substitute for raw coordinates for conductivity and pH. The block-bootstrap intervals on discrimination overlap between the feature sets compared, so the ablation does not separate them (\nameref{S5_Text}); \nameref{S6_Fig} shows one illustrative surface, read as an ordering and not as a probability. Its stack comprises the national aquifer classification~\cite{CGWB-ManualAquiferMapping}, SoilGrids 2.0 soil properties~\cite{Poggio2021}, ESA WorldCover 10 m 2021 v200 land cover~\cite{Zanaga2022} and WorldPop population density~\cite{WorldPop2020}.

\subsection*{External transfer assessment}

Transfer was assessed per archive \texttimes{} parameter pair. For each parameter a separate Model A fit was trained on GEMStat records before a cut year, on 33 features that add nearest-station neighbour features and month to the 24 above (\nameref{S1_Text}). These transfer fits are not the seven validation fits and were not evaluated under the three designs, so the optimism measured under the designs is not carried to the transferred scores. The fixed fits were scored on the external archives with no tuning, threshold selection or recalibration after the external data were seen, once every external site sharing a 0.1\textdegree{} cell with a training station had been discarded. Of 24 possible pairs, 13 were evaluated and 11 are reported as untested rather than negative.

The verdict comes from three interval tests rather than from a threshold on discrimination, each passed when a 95\% bootstrap interval lies wholly above its null value: (i) the site-level rank correlation between predicted risk and observed severity, where severity is log\textsubscript{10} concentration or, for pH, the distance from the centre of 6.5--8.5 (site bootstrap; above 0); (ii) the rank correlation within 1\textdegree{} blocks (block bootstrap; above 0); and (iii) lift in the top predicted decile (site bootstrap; above 1). With fewer than two tests evaluable a pair has \textit{insufficient support}, a gap in the evidence that is never merged with a negative result; if every evaluable test passes, \textit{rank transfer is supported}; if none passes, it is \textit{not supported}; otherwise it is \textit{partial or scale-dependent}. Only the outcome of each test was retained, not its statistic or interval, so a verdict carries no stated sampling uncertainty. For every pair, the retained results also give two quantities that are not verdict tests, each with a site-clustered bootstrap interval: threshold discrimination (threshold ROC-AUC) and a detected-record severity correlation, the Spearman rank correlation between predicted risk and observed severity among detected records only. They also give the calibration slope, Brier skill against a reference forecast equal to the external base rate, and confusion counts at a 0.5 cut-point (\nameref{S1_Text}). The retained calibration intercept was estimated jointly with the slope and so is not calibration-in-the-large, and no criterion for adequate calibration was fixed in advance. A supported verdict authorises prioritised confirmatory sampling only, at the level of observations within the archive.

The Kenyan and Malawian archive summaries were recomputed from the raw archives for this article. The Indian state-level results are the outputs of the original analysis, checked for internal consistency but not re-run, and the validation, transfer and ablation results were checked against retained result exports rather than re-run for this article (\nameref{S1_Text}).

\subsection*{Technology evidence and action rules}

The exact-product catalogue comprises 23 named commercial products assessed against 54 criteria, a 1,242-cell grid in which each cell records its value, source, locus and evidence class (\nameref{S4_Table}). Not reported, not researched, conflicting and range-only evidence are kept distinct, and evidence about an adjacent product cannot verify a named model. No record states how the 23 products were chosen or which sources were searched and when, so the product counts describe the documentary record as located. A criterion meets the documentation threshold when it is documented for at least 12 of the 23 products, counted on its primary database field (strict) or on any database field mapped to it (any admissible). The sufficiency rule requires all 11 active criteria (\nameref{S1_Text}) to be documented; threshold-near classification error, invalid-test rate and testing frequency are not among the active criteria.

Conditional comparison of the ten technology archetypes used the same 11 criteria under seven declared weighting scenarios and 13 cost mappings, with 40,000 Monte Carlo draws per scenario from a Dirichlet distribution centred on the scenario's weights (\nameref{S1_Text}, \nameref{S7_Table}). The technique for order of preference by similarity to ideal solution (TOPSIS) supplied the primary rank~\cite{Chakraborty2022}, and stochastic multicriteria acceptability analysis (SMAA) represented weight uncertainty as a distribution over the feasible weights~\cite{TervonenLahdelma2007}. The comparison is not screened for the analytes a scenario targets, and no practitioner weights were elicited.

For each monitoring context a rule returns an action (\nameref{S2_Table}): for Kenya and Malawi from the data status alone, whatever the exceedance count, and for India from a per-parameter evidence state that we assigned from the transfer verdict, record support and identity limits (\nameref{S1_Text} gives both branches). The evidence states and the rule were recorded on 13 August 2026, after the transfer verdicts had been computed, and no dated protocol records when the action level, support requirement, verdict rule, documentation threshold or active criteria were fixed. The product-level decision is a single determination over the 23-product catalogue, applied to all 26 contexts, not 26 independent determinations.

\subsection*{Use of artificial-intelligence tools}

OpenAI ChatGPT (version not recorded) was used to set up the project's initial prompts and scripts, and Anthropic Claude Opus 5, through Claude Code, to write and run data-processing, verification, table and figure scripts; to search for and select literature, record the passage supporting each citation and assemble the reference list; to draft and revise every section of the text and the supporting information; and to review drafts, with accepted recommendations applied. The validation and transfer model runs were not produced with these tools, and no tool altered or imputed a measurement. Table and figure values were re-derived by script from source files, every number in the text was checked against a register of sourced values, and each cited sentence was checked against the passage recorded from its opened source. \nameref{S10_Text} lists each activity and its check; at the time of writing, the authors' review of the text and citation records was outstanding.

\FloatBarrier

\section*{Results}

The applications run from India, where the accepted archives are largest, to the local archives of Kwale County in Kenya and of Malawi, and then to published evidence from Bangladesh and Ghana. Across the 26 monitoring contexts the rule returned eight distinct actions, each with the evidence required to reassess it (Fig~\ref{fig:action-distribution}, Table~\ref{table1}). The most frequent were \textit{local signal plus direct confirmation}, in 9 contexts, and \textit{monitoring baseline required}, in 5; six contexts hold no records, and their actions follow from that absence. An action for every context reflects the design of the rules, not their accuracy.\label{callout:fig:action-distribution}\label{callout:table1}

\begin{figure}[!ht]
\caption{\textbf{Actions returned in the 26 monitoring contexts.} Number of contexts, each a country \texttimes{} source archive \texttimes{} parameter combination, returning each of the eight actions, by country (India 8, Kenya 8, Malawi 10). Data: \nameref{S9_Data}.}
\label{fig:action-distribution}
\end{figure}

\nolinenumbers
\begin{spacing}{1}
{\footnotesize\setlength{\tabcolsep}{3pt}\setlength{\LTcapwidth}{\linewidth}
\begin{longtable}{>{\raggedright\arraybackslash}p{0.2385\linewidth}>{\raggedleft\arraybackslash}p{0.0769\linewidth}>{\raggedright\arraybackslash}p{0.2077\linewidth}>{\raggedright\arraybackslash}p{0.4262\linewidth}}
\caption{\textbf{Monitoring actions and reassessment evidence across 26 contexts.}}\label{table1} \\
\hline
\textbf{Monitoring context} & \textbf{Records (n)} & \textbf{Action} & \textbf{Evidence required to reassess the action} \\ \thickhline
\endfirsthead
\multicolumn{4}{@{}l}{\footnotesize\emph{Table 1. (continued)}} \\
\hline
\textbf{Monitoring context} & \textbf{Records (n)} & \textbf{Action} & \textbf{Evidence required to reassess the action} \\ \thickhline
\endhead
India, CGWB Year Book: fluoride & 35,843 & Rank prioritisation plus direct confirmation & Repair longitudinal well identity and run same-sample field-versus-reference validation. \\ \hline
India, CGWB Year Book: nitrate & 34,642 & Rank prioritisation plus direct confirmation & Repair persistent well/date/sample identity and collect repeat measurements before any temporal model. \\ \hline
India, CGWB portal: arsenic & 9,927 & Targeted direct reference measurement & Recover method and detection-limit metadata and build a defensible detected/non-detected analytical population. \\ \hline
India, CGWB portal: conductivity & 173,029 & Rank prioritisation plus direct confirmation & Governance review of the 1,500 \textmu{}S/cm project convention plus prospective sentinel-versus-spot evaluation; major-ion chemistry at flagged wells. \\ \hline
India, CGWB portal: pH & 162,729 & Direct measurement required & Collect representative local labels at sentinel sites and assess whether a separate local recalibration or model is useful; pH departures reported by direction. \\ \hline
India, CGWB portal: TDS & 30,151 & Rank prioritisation plus direct confirmation & Prospective paired conductivity/TDS sampling across relevant water matrices and seasons; recovery of how the portal TDS values were derived. \\ \hline
India, no accepted archive: \textit{E.~coli} & -- & Monitoring baseline required & Acquire traceable India microbial measurements with household/source context and compatible denominators. \\ \hline
India, no accepted archive: turbidity & -- & Monitoring baseline required & Acquire a machine-readable India archive with site, date, method, unit and valid turbidity results. \\ \hline
Kenya, Kwale groundwater samples: arsenic & 140 & Local signal plus direct confirmation & Repeat same-site arsenic measurement across additional seasons and confirm threshold classifications with a traceable reference workflow. \\ \hline
Kenya, Kwale groundwater samples: conductivity & 140 & Local signal plus direct confirmation & Repeat same-site conductivity measurement across additional seasons and confirm threshold classifications with a traceable reference workflow. \\ \hline
Kenya, Kwale groundwater samples: \textit{E.~coli} & -- & Monitoring baseline required & Acquire traceable local \textit{E.~coli} records with site, date, unit, method and quality-control fields. \\ \hline
Kenya, Kwale groundwater samples: fluoride & 139 & Local signal plus direct confirmation & Repeat same-site fluoride measurement across additional seasons and confirm threshold classifications with a traceable reference workflow. \\ \hline
Kenya, Kwale groundwater samples: nitrate & 139 & Local signal plus direct confirmation & Repeat same-site nitrate measurement across additional seasons and confirm threshold classifications with a traceable reference workflow. \\ \hline
Kenya, Kwale groundwater samples: pH & 140 & Local signal plus direct confirmation & Repeat same-site pH measurement across additional seasons and confirm threshold classifications with a traceable reference workflow. \\ \hline
Kenya, Kwale groundwater samples: TDS & -- & Monitoring baseline required & Acquire traceable local TDS records with site, date, unit, method and quality-control fields. \\ \hline
Kenya, Kwale groundwater samples: turbidity & -- & Monitoring baseline required & Acquire traceable local turbidity records with site, date, unit, method and quality-control fields. \\ \hline
Malawi, laboratory snapshot: arsenic & 10 & Data repair required & Recover arsenic units, method, LOD/LOQ and non-detect semantics before threshold analysis. \\ \hline
Malawi, laboratory snapshot: conductivity & 31 & Data repair required & Recover the conductivity scale/unit convention and raw instrument export; do not repair by assumption. \\ \hline
Malawi, laboratory snapshot: \textit{E.~coli} & -- & Organism-specific baseline required & Acquire organism-specific \textit{E.~coli} measurements with sample volume, reference method and response pathway. \\ \hline
Malawi, laboratory snapshot: faecal coliform & 19 & Organism-specific baseline required & Acquire organism-specific \textit{E.~coli} measurements with sample volume, reference method and response pathway. \\ \hline
Malawi, laboratory snapshot: fluoride & 2 & Data repair required & Recover the fluoride missing-value codebook and valid observations before threshold analysis. \\ \hline
Malawi, laboratory snapshot: nitrate & 31 & Local signal plus direct confirmation & Expand repeat-site nitrate coverage and retain method, calibration, operator and quality-control metadata. \\ \hline
Malawi, laboratory snapshot: pH & 31 & Local signal plus direct confirmation & Expand repeat-site pH coverage and retain method, calibration, operator and quality-control metadata. \\ \hline
Malawi, laboratory snapshot: TDS & 31 & Local signal plus direct confirmation & Expand repeat-site TDS coverage and retain method, calibration, operator and quality-control metadata. \\ \hline
Malawi, laboratory snapshot: turbidity & 22 & Local signal plus direct confirmation & Expand repeat-site turbidity coverage and retain method, calibration, operator and quality-control metadata. \\ \hline
Malawi, logger series: conductivity & 714 & High-frequency sentinel plus chemical confirmation & Pair the two persistently high-conductivity sites and comparison sites with direct chemistry, calibration, fouling, downtime and energy records. \\ \hline
\end{longtable}
\begin{flushleft}
1. The framework names no product in any of the 26 contexts (one catalogue-wide decision), and no row is a national estimate.

2. An en dash means the count is not available (no accepted measurement population), not zero exceedances; absence of records is not evidence of safety.

3. Counts are accepted records, except held rows that failed acceptance (Malawi conductivity, fluoride and arsenic) and logger site-days; Malawi TDS records share the mixed scales that led to data repair for conductivity (Results).

4. The last column names the evidence that would change an action, not the action it would change to, who supplies it, or its sample size, frequency, method or cost. Full matrix and action definitions: \nameref{S2_Table}. LOD/LOQ, limits of detection and quantification.
\end{flushleft}
}
\end{spacing}
\linenumbers

\FloatBarrier

\subsection*{Decision-relevant cells and record support in India}

The Indian archives resolve to 186 state or Union Territory \texttimes{} parameter cells in 31 states and Union Territories, of which 160 met the support requirement for description; accepted record counts range from 173,029 for conductivity to 9,927 for arsenic (Table~\ref{table2}).\label{callout:table2}

\begin{table}[!ht]
\centering
\caption{\textbf{Indian monitoring records and decision-relevant cells by parameter.}}
{\footnotesize\setlength{\tabcolsep}{3pt}
\begin{tabular}{>{\raggedright\arraybackslash}p{0.1077\linewidth}>{\raggedleft\arraybackslash}p{0.0846\linewidth}>{\raggedleft\arraybackslash}p{0.0923\linewidth}>{\raggedleft\arraybackslash}p{0.0692\linewidth}>{\raggedright\arraybackslash}p{0.1077\linewidth}>{\raggedright\arraybackslash}p{0.1769\linewidth}>{\raggedleft\arraybackslash}p{0.0923\linewidth}>{\raggedleft\arraybackslash}p{0.0800\linewidth}>{\raggedleft\arraybackslash}p{0.0738\linewidth}}
\hline
\textbf{Parameter} & \textbf{Accepted records (n)} & \textbf{Records above threshold (n)} & \textbf{Sites (n)} & \textbf{Exceedance threshold} & \textbf{Record fraction [95\% Wilson interval]} & \textbf{Described cells (of 31)} & \textbf{Decision-relevant cells (n)} & \textbf{Cells with LOSO support (of 31)} \\ \thickhline
Conductivity & 173,029 & 39,115 & 24,128 & \textgreater{}1,500 \textmu{}S/cm (operational convention) & 0.2261 [0.2241, 0.2280] & 30 & 1 & 22 \\ \hline
pH & 162,729 & 10,409 & 23,704 & outside 6.5--8.5 pH units & 0.0640 [0.0628, 0.0652] & 30 & 4 & 21 \\ \hline
Fluoride & 35,843 & 2,344 & 32,561 & \textgreater{}1.5 mg/L & 0.0654 [0.0629, 0.0680] & 27 & 6 & 14 \\ \hline
Nitrate & 34,642 & 6,403 & 32,206 & \textgreater{}50 mg/L as NO\textsubscript{3} & 0.1848 [0.1808, 0.1890] & 27 & 2 & 18 \\ \hline
TDS & 30,151 & 5,442 & 14,632 & \textgreater{}1,000 mg/L & 0.1805 [0.1762, 0.1849] & 27 & 4 & 14 \\ \hline
Arsenic & 9,927 & 31 & 7,931 & \textgreater{}0.010 mg/L & 0.0031 [0.0022, 0.0044] & 19 & 0 & 0 \\ \hline
\end{tabular}
\begin{flushleft}
1. Record fractions carry Wilson 95\% intervals, which treat records as independent. A cell is decision-relevant when its interval straddles the assumed 0.10 action level (17 of 160 described cells, a lower bound and not a ranking of severity). LOSO, leave-one-state-out.

2. Conductivity uses 1,500~\textmu{}S/cm, an operational project convention, not a health guideline. Nitrate is classified at 50~mg/L as NO\textsubscript{3}; the transfer analysis uses 10~mg/L as N, the United States maximum contaminant level.

3. Fluoride and nitrate come from the CGWB Year Book, the other parameters from the CGWB portal; records cover all 31 states and Union Territories except arsenic (23 of them). Per-cell values, support and revisit cadence: \nameref{S3_Table}.
\end{flushleft}
}
\label{table2}
\end{table}

At an assumed action level of 0.10, 17 of the 160 eligible cells are decision-relevant, in 15 states and Union Territories (Fig~\ref{fig:india-map}): 6 fluoride cells, 4 pH, 4 TDS, 2 nitrate and 1 conductivity. None of the 19 eligible arsenic cells is decision-relevant; the arsenic record fraction across all accepted records, 0.0031 (Table~\ref{table2}), is not interpretable as a monitoring signal until units, method and detection limits are recovered. These 17 cells are where further records could change a state-level classification; they are not a ranking of severity, prevalence or exposure, and no action in Table~\ref{table1} depends on them.\label{callout:fig:india-map}

\begin{figure}[!ht]
\caption{\textbf{Decision-relevant cells in the Indian archives.} Markers locate the 17 state or Union Territory \texttimes{} parameter cells whose Wilson 95\% intervals on the exceedance proportion straddle the assumed 0.10 action level, numbered from the widest interval to the narrowest (an order of uncertainty, not severity). Markers are administrative anchors, not sampling sites. Boundaries: Natural Earth. Data: \nameref{S9_Data}.}
\label{fig:india-map}
\end{figure}

Each of the 17 Wilson intervals straddles the action level (Fig~\ref{fig:india-intervals}); the widest, for pH in Nagaland, rests on 44 accepted records and the narrowest, for conductivity in Uttar Pradesh, on 7,280. The count depends on the level, with 29 cells in 16 states and Union Territories at 0.05 and 4 cells in 4 at 0.25. The intervals treat records as independent although sites are revisited, so the counts are lower bounds, and a cell resolved below the action level is not one where monitoring may stop.\label{callout:fig:india-intervals}

\begin{figure}[!ht]
\caption{\textbf{Exceedance proportions in the 17 decision-relevant cells.} Points and bars show record proportions with Wilson 95\% intervals; n is the number of accepted records. Numbering follows Fig~\ref{fig:india-map} and the dashed line marks the assumed 0.10 action level. The intervals treat records as independent, so 17 is a lower bound. TDS, total dissolved solids. Data: \nameref{S9_Data}.}
\label{fig:india-intervals}
\end{figure}

The other 143 eligible cells are resolved, with intervals wholly above or wholly below the action level. Rajasthan shows the upper side: of 12,349 conductivity records at 1,641 sites between 2000 and 2024, 6,676 exceeded the 1,500~\textmu{}S/cm project convention (54.1\%, Wilson 95\% interval 53.2\% to 54.9\%), and of 3,131 TDS records at 1,085 sites, 1,627 exceeded 1,000~mg/L (52.0\%, 50.2\% to 53.7\%). No Rajasthan conductivity or TDS record states its analytical method.

Record support and revisit cadence set what further analysis the archives can carry (Fig~\ref{fig:india-support-cadence}). A leave-one-state-out evaluation would be supported in 89 of the 186 cells, including 9 of the 17 decision-relevant cells, and in none of the 31 arsenic cells. In all 145 cells with a recorded median revisit interval the median is annual or longer; 48 are exactly 365 days and the longest is 1,534 days. Because 46 of the 145 medians come from Year Book tables, whose intervals cannot fall below the spacing between editions, and some rest on a single interval, the cadence result describes the archive as held rather than the programme as now designed. No action in Table~\ref{table1} prescribes a sampling frequency.\label{callout:fig:india-support-cadence}

\begin{figure}[!ht]
\caption{\textbf{Record support and revisit cadence in India.} (A) Cells per parameter, of 31, supporting description (grey bar), a leave-one-state-out evaluation (hatched bar) and decision-relevant cells (diamond): 160, 89 and 17 of 186 in all. (B) Median revisit interval in the 145 cells with one recorded; Year Book cells cannot fall below the spacing between editions. TDS, total dissolved solids. Data: \nameref{S9_Data}.}
\label{fig:india-support-cadence}
\end{figure}

\FloatBarrier

\subsection*{Model performance, transfer and the Indian actions}

Two model questions follow. The first asks how far a performance estimate depends on validation design; it is answered on the multinational fitting archive and sets no action, but it fixes which design a claimed use requires. The second asks whether ordering and probability carry over to archives that took no part in fitting, and it bears on the Indian actions. The two questions use separate fits: the validation-design fits below are not the transfer fits. Validation design changed the performance estimate although the model specification and data were the same, with the model refitted on each design's training partition (Fig~\ref{fig:design-discrimination}, Table~\ref{table3}). Across seven tasks from the multinational fitting archive, mean test-partition ROC-AUC was 0.9396 under random-row splitting, 0.8688 under station-grouped splitting and 0.7696 under 3\textdegree{} block splitting. The optimism, random-row minus block ROC-AUC, averaged 0.1701, computed from unrounded means, and ranged from 0.0768 for fluoride to 0.2927 for nitrate. The optimism estimate comes from one partition per design on seven tasks from one predominantly surface-water archive, with no interval.\label{callout:fig:design-discrimination}\label{callout:table3}

\begin{figure}[!ht]
\caption{\textbf{Model discrimination under three validation designs.} Test-partition ROC-AUC (area under the receiver operating characteristic curve) for seven tasks under random-row, station-grouped and spatial-block splitting, one partition per design. Optimism is random-row minus spatial-block ROC-AUC, computed from unrounded values. Data: Table~\ref{table3} and \nameref{S9_Data}.}
\label{fig:design-discrimination}
\end{figure}

\begin{table}[!ht]
\centering
\caption{\textbf{Model performance by validation design and information channel.}}
{\footnotesize\setlength{\tabcolsep}{3pt}
\begin{tabular}{>{\raggedright\arraybackslash}p{0.1015\linewidth}>{\raggedleft\arraybackslash}p{0.0954\linewidth}>{\raggedleft\arraybackslash}p{0.0677\linewidth}>{\raggedleft\arraybackslash}p{0.0754\linewidth}>{\raggedleft\arraybackslash}p{0.0677\linewidth}>{\raggedleft\arraybackslash}p{0.0923\linewidth}>{\raggedleft\arraybackslash}p{0.0677\linewidth}>{\raggedleft\arraybackslash}p{0.0923\linewidth}>{\raggedleft\arraybackslash}p{0.0662\linewidth}>{\raggedleft\arraybackslash}p{0.0662\linewidth}>{\raggedleft\arraybackslash}p{0.0662\linewidth}}
\hline
\textbf{Parameter} & \textbf{Prevalence} & \textbf{ROC-AUC random row} & \textbf{ROC-AUC station grouped} & \textbf{ROC-AUC spatial block} & \textbf{Optimism} & \textbf{Spatial PR-AUC} & \textbf{Spatial calibration slope} & \textbf{Model A ROC-AUC} & \textbf{Model B ROC-AUC} & \textbf{Model B \textminus{} A} \\ \thickhline
Arsenic & 0.0833 & 0.9626 & 0.9230 & 0.8160 & 0.1466 & 0.2710 & 0.8922 & 0.9302 & 0.9530 & +0.0228 \\ \hline
Nitrate & 0.0263 & 0.9727 & 0.8527 & 0.6800 & 0.2927 & 0.0743 & 0.8363 & 0.9444 & 0.9779 & +0.0335 \\ \hline
Conductivity & 0.0416 & 0.9854 & 0.9411 & 0.8596 & 0.1258 & 0.3428 & 0.7001 & 0.9633 & 0.9858 & +0.0225 \\ \hline
Suspended solids & 0.2277 & 0.8683 & 0.8203 & 0.7463 & 0.1221 & 0.4633 & 0.8541 & 0.8508 & 0.8724 & +0.0215 \\ \hline
Turbidity & 0.8323 & 0.9293 & 0.8615 & 0.7440 & 0.1853 & 0.9272 & 0.6749 & 0.9121 & 0.9426 & +0.0305 \\ \hline
Fluoride & 0.0228 & 0.9830 & 0.9336 & 0.9062 & 0.0768 & 0.1912 & 0.1925 & -- & -- & -- \\ \hline
pH & 0.0715 & 0.8760 & 0.7496 & 0.6348 & 0.2412 & 0.1855 & 0.2979 & 0.8298 & 0.8940 & +0.0643 \\ \hline
\textbf{Mean of 7} & n/a & \textbf{0.9396} & \textbf{0.8688} & \textbf{0.7696} & \textbf{0.1701} & n/a & n/a & n/a & n/a & n/a \\ \hline
\end{tabular}
\begin{flushleft}
1. Optimism is random-row minus spatial-block ROC-AUC, computed before rounding. Only spatial-block splitting speaks to held-out areas within the sampled geography, and no design addresses transfer to an unmonitored country. Prevalence is a task's share of records above threshold, the reference for PR-AUC.

2. Model A (33 features, separate from the validation-design fits) and Model B (12 site-history features) were compared in forward-chained yearly folds on records with at least two earlier same-site measurements; calibration slopes for both models are in \nameref{S1_Text}.

3. An en dash means not available (fluoride has no paired comparison); n/a means not applicable. DeLong p-values are omitted: the exported 0.00e+00 is below display precision, and the test assumes independent observations, which revisited stations violate. Values are from retained result tables (\nameref{S1_Text}).
\end{flushleft}
}
\label{table3}
\end{table}

Under block splitting, precision--recall AUC exceeded each task's overall prevalence for all seven tasks, from 0.0743 for nitrate to 0.9272 for turbidity (Fig~\ref{fig:prevalence-prauc}). Turbidity is the one task in which exceedance is the majority class, with a prevalence of 0.8323; the strict no-skill baseline, the prevalence of the scored test partition, is not reported in the retained tables.\label{callout:fig:prevalence-prauc}

\begin{figure}[!ht]
\caption{\textbf{Spatial-block PR-AUC against task prevalence.} Precision--recall AUC (PR-AUC) under spatial-block splitting for the seven tasks, beside each task's overall prevalence, the share of its records above threshold. The test-partition prevalence, the strict no-skill baseline, was not retained. Data: \nameref{S9_Data}.}
\label{fig:prevalence-prauc}
\end{figure}

In forward-chained yearly folds, Model B exceeded Model A on each of the six parameters with a paired comparison, by 0.0215 to 0.0643 in ROC-AUC (Fig~\ref{fig:channel-comparison}). The advantage needs earlier measurements at the same site, so it is unavailable where monitoring is absent, and it is not evidence of spatial generalisation.\label{callout:fig:channel-comparison}

\begin{figure}[!ht]
\caption{\textbf{Model A and Model B in forward-chained yearly folds.} ROC-AUC of Model A (environmental, positional and neighbour features) and Model B (earlier same-site measurements) for the six tasks with a paired comparison; each year is predicted from earlier years, on records with at least two earlier same-site measurements. Data: \nameref{S9_Data}.}
\label{fig:channel-comparison}
\end{figure}

Ordering transferred for five of the 13 archive \texttimes{} parameter pairs evaluated, six of them Indian (Fig~\ref{fig:transfer-verdicts}): conductivity and TDS in the CGWB portal, fluoride and nitrate in the CGWB Year Book, and TDS in WQP. TDS was supported in both archives in which it was tested (threshold ROC-AUC 0.7081 in the CGWB portal and 0.7406 in WQP), and pH in neither (0.4870 and 0.4886; the CGWB portal pH interval, 0.4784 to 0.4964, lies wholly below 0.5). Arsenic in the CGWB portal had insufficient support, with one evaluable test on 119 scored observations at 116 sites, and its ordering was inverted: ROC-AUC 0.3170 (95\% interval 0.2043 to 0.4450). The nitrate ordering was evaluated at 10~mg/L as N, and its use against 50~mg/L as NO\textsubscript{3} assumes that the ordering does not depend on that difference. Scored observations are fewer than accepted records for every Indian pair; the site exclusion accounts for part of the difference, and the rest is not recorded.\label{callout:fig:transfer-verdicts}

\begin{figure}[!ht]
\caption{\textbf{External rank-transfer verdicts by archive and parameter.} Cells show the verdict, threshold ROC-AUC and scored observations; 13 of 24 pairs were evaluated. Verdicts follow three interval tests (site-level and within-block rank correlation, and lift in the top predicted decile), not a ROC-AUC cut-off; insufficient support is not a negative result. CGWB, Central Ground Water Board; TDS, total dissolved solids. Data: \nameref{S9_Data}.}
\label{fig:transfer-verdicts}
\end{figure}

The detected-record severity correlation, which is not a verdict test, had a 95\% interval wholly above 0 for 9 of the 13 pairs, including four without supported transfer. Among the five supported pairs it ranged from 0.2025 for nitrate in the CGWB Year Book to 0.4778 for TDS in WQP, and their threshold ROC-AUC ranged from 0.6107 to 0.7426 (Fig~\ref{fig:transfer-ordering}); the correlation describes detected records only and enters no verdict.\label{callout:fig:transfer-ordering}

\begin{figure}[!ht]
\caption{\textbf{Discrimination and severity ordering under external transfer.} (A) Threshold ROC-AUC and (B) detected-record severity correlation (Spearman, predicted risk against observed severity among detected records), each with a site-clustered 95\% bootstrap interval, for the 13 evaluated pairs. Severity is log\textsubscript{10} concentration, or for pH the distance from the centre of 6.5--8.5; the correlation is not a verdict test. Marker shape and colour give the verdict. Data: \nameref{S9_Data}.}
\label{fig:transfer-ordering}
\end{figure}

Probability did not transfer with ordering (Fig~\ref{fig:transfer-calibration}). The calibration slope was below 1 on all 13 evaluated pairs, between 0.2383 and 0.4017 on the five pairs whose ordering transferred, and at or below 0 on 4; within the fitting archive the block-design slopes were already below 1 for all seven tasks, between 0.1925 and 0.8922. Brier skill was negative on all six Indian pairs and on 11 of the 13 overall; the two non-negative values, 0.0470 for \textit{E.~coli} and 0.0169 for arsenic in WQP, belong to pairs without supported rank transfer. The conductivity base rate was 0.0638 in the transfer training data, 0.2261 among all accepted conductivity records and 0.2364 among the scored CGWB portal observations, a shift that any probability carried from the fitting archive would have to absorb. Calibration-in-the-large was not assessed and local recalibration was not run, so no transferred score can be read as a local exceedance probability. At a fixed 0.5 cut-point most exceeding observations were not flagged; for conductivity in the CGWB portal, sensitivity was 0.0787 (\nameref{S1_Text} gives the counts).\label{callout:fig:transfer-calibration}

\begin{figure}[!ht]
\caption{\textbf{Probability calibration under external transfer.} (A) Calibration slope and (B) Brier skill against a reference forecast equal to the external base rate, for the 13 evaluated pairs; reference lines mark slope 1 and skill 0. No local recalibration was run. Marker shape and colour give the verdict; nitrate is scored at 10~mg/L as N. Data: \nameref{S9_Data}.}
\label{fig:transfer-calibration}
\end{figure}

The evidence states we assigned from the transfer verdicts, record support and identity limits returned the Indian actions (Table~\ref{table1}). Conductivity, TDS, nitrate and fluoride returned \textit{rank prioritisation plus direct confirmation}: ordering may set the order in which sources are confirmed, every flagged source still needs a direct measurement, low-ranked sources stay under routine measurement, and no ranking is released with this article. pH returned \textit{direct measurement required}. Arsenic returned \textit{targeted direct reference measurement}, pending method and detection-limit metadata; this study sets no basis for choosing the targets and did not assess laboratory access or cost. Turbidity and \textit{E.~coli}, with no accepted Indian measurement population in the archives analysed, returned \textit{monitoring baseline required}; CGWB lists microbial contaminants among additional parameters to be monitored~\cite{CGWB-SOP-GWQDataAnalysis}, so the evidence may lie with drinking-water agencies rather than in a new survey. Tables~\ref{table1} and~\ref{table4} give the evidence that would change each action.

\subsection*{Repeat measurement in Kwale County, Kenya}

The Kwale County archive describes groundwater and springs sampled in March and June 2016 and is used without any model. It records 81 sampling locations (Fig~\ref{fig:kwale-map}) across 71 named localities; once surface-water samples are excluded, 73 locations remain, with 139 or 140 records per parameter. Locations are counted as distinct recorded coordinates, so the count depends on how coordinates were recorded. Measured pH lay outside 6.5--8.5 in 41 of 140 records (29.3\%, Wilson 95\% interval 22.4\% to 37.3\%), all of them below 6.5, and conductivity exceeded the 1,500~\textmu{}S/cm convention in 19 of 140 (13.6\%, 8.9\% to 20.2\%). Nitrate exceeded 50~mg/L as NO\textsubscript{3} in 5 of 139 (3.6\%, 1.5\% to 8.1\%), fluoride exceeded 1.5~mg/L in 3 of 139 (2.2\%, 0.7\% to 6.2\%) and arsenic exceeded 10~\textmu{}g/L in 3 of 140 (2.1\%, 0.7\% to 6.1\%). Most locations were sampled twice, so these intervals understate the uncertainty.\label{callout:fig:kwale-map}

\begin{figure}[!ht]
\caption{\textbf{Sampling locations of the Kwale County archive, Kenya.} The 81 sampling locations of the Kwale County archive, from two 2016 campaigns~\cite{OWD-grdwtrsmpkwale}, reprojected from UTM zone 37S on the Arc 1960 datum without a datum shift; no coastline is drawn at this scale. The text's 73 locations exclude surface water. Inset: Kenya. Openwashdata package under CC BY 4.0. Data: \nameref{S9_Data}.}
\label{fig:kwale-map}
\end{figure}

Two of the three fluoride exceedances, and two of the 19 conductivity exceedances, come from one thermal spring with a recorded temperature of about 59 \textdegree{}C, and the third fluoride exceedance and one June conductivity crossing are unverified values (\nameref{S1_Text}). The counts include these records as the archive holds them.

At the 67 sites sampled in both campaigns (66 with nitrate and fluoride results in both), 15 of the 24 sites with pH outside the range in March were still outside it in June, and 2 more were outside it in June only, so pH classification changed at 11 of the 67 sites. Of the 10 sites above the conductivity convention in March, 8 remained above it and 1 more was above it in June only. For nitrate, fluoride and arsenic, one site was above the threshold in both campaigns, and 2, 1 and 1 sites respectively were above it in June only (\nameref{S11_Table}); the fifth nitrate exceedance is at a location sampled once. The source labels the campaigns dry-season and wet-season, but one campaign in each season of a single year does not describe a seasonal cycle.

For pH, conductivity, nitrate, fluoride and arsenic, the rule returned \textit{local signal plus direct confirmation} from the data status alone: repeat measurement at the same sites across further seasons, with threshold classifications confirmed by a reference method the rule does not name (Table~\ref{table1}). Contexts that share the archive and the action could be met by one sampling round. The action sets what to measure next. It does not defer confirmation or a programme's response where a source was recorded above a WHO health-based guideline value in 2016. TDS, turbidity and \textit{E.~coli}, for which the archive holds no records, returned \textit{monitoring baseline required}.

\subsection*{Data repair and sentinel monitoring in Malawi}

The Malawi laboratory snapshot holds one record per water point for 32 water points sampled between December 2017 and July 2019 (Fig~\ref{fig:malawi-map}). Measured pH lay outside 6.5--8.5 in 7 of 31 records (22.6\%, Wilson 95\% interval 11.4\% to 39.8\%), all between 6.24 and 6.42, and the archive flags each as within the Malawian standard it applies. No nitrate record of 31 exceeded 50~mg/L as NO\textsubscript{3}, which leaves an upper Wilson bound of 11.0\%. Turbidity exceeded the 1~NTU rule in 2 of 22 records (9.1\%, 2.5\% to 27.8\%), both flagged as within the Malawian standard. Faecal coliform was detectable in 4 of 19 samples, all collected in February 2019 and none with a recorded analysis date. Faecal coliform is not \textit{E.~coli}, but WHO treats thermotolerant coliforms as an acceptable alternative indicator of faecal pollution~\cite{WHO2022GDWQ}, so detection at those four water points would ordinarily warrant inspection and resampling whatever the organism-specific evidence gap.\label{callout:fig:malawi-map}

\begin{figure}[!ht]
\caption{\textbf{Malawian laboratory water points and conductivity logger sites.} The laboratory snapshot of 32 Malawian water points~\cite{OWD-boreholelabdata} and the conductivity logger series at 19 sites~\cite{OWD-mwgroundwaterdata}. Nearby markers overlap at this scale. Boundaries: Natural Earth; openwashdata packages under CC BY 4.0. Data: \nameref{S9_Data}.}
\label{fig:malawi-map}
\end{figure}

Four parameters could not be read as recorded: conductivity failed a unit-consistency check, TDS is recorded on mixed scales in 10 of 31 records, so no TDS count is reported, fluoride held only a placeholder code, and arsenic had no documented unit. Conductivity, fluoride and arsenic returned \textit{data repair required}, to be made from recovered records and not by assumption; only the originating laboratory or the data publisher can make it, and the rule has no fallback if the records cannot be recovered. The laboratory pH, TDS, nitrate and turbidity contexts returned \textit{local signal plus direct confirmation}. For nitrate, with no exceedance in 31 records, the action amounts to wider coverage; the TDS records passed acceptance as recorded and so returned local signal plus direct confirmation, but they share the mixed scales that returned data repair for conductivity, and applying the conductivity check to them would also return data repair. \textit{E.~coli}, with no organism-specific measurement, and the faecal coliform context both returned \textit{organism-specific baseline required} (Table~\ref{table1}).

The logger extract holds conductivity from 19 sites in 9 districts, reduced to 714 daily site medians between 31 December 2024 and 7 February 2025, a window of 39 days within Malawi's wet season from November to April~\cite{BGS-Earthwise-Malawi}. On 76 of the 714 site-days, conductivity exceeded the 1,500~\textmu{}S/cm convention: two sites, both in Chikwawa District, exceeded it on every recorded day and no other site on any day, so across 695 successive site-day transitions no site changed state. The two sites have median conductivities of about 5,690 and 2,160~\textmu{}S/cm, and the next highest site median is about 1,360~\textmu{}S/cm. The series cannot identify the cause: conductivity is not specific to any contaminant, the logger does not record whether values are temperature-compensated, and the extract covers no dry-season days. The rule returned \textit{high-frequency sentinel plus chemical confirmation}, pairing the two persistent sites, and comparison sites, with direct chemistry and with instrument calibration, fouling, downtime and energy records; the analytes, comparison sites, chemistry frequency and logging period are not set by this study.

\subsection*{Screening and testing-service design in Bangladesh and Ghana}

In Bangladesh, where no archive was analysed and no model transferred, published kit evaluations bear on whether a field kit can classify a well near the arsenic threshold accurately enough to leave it in service. In a blind evaluation of 123 wells against laboratory measurement, with 65 tested twice, one kit classified 110 wells correctly against the WHO guideline of 10~\textmu{}g/L and 113 against the Bangladesh standard of 50~\textmu{}g/L~\cite{George2012}. All 13 of its misclassifications against 10~\textmu{}g/L were underestimates, in wells between 10 and 27~\textmu{}g/L~\cite{George2012}. Against 50~\textmu{}g/L its over- and underestimates were evenly distributed, and the evaluation omitted the kit's hydrogen-sulphide oxidation step~\cite{George2012}. A comparison of eight commercial kits with hydride-generation atomic absorption spectrometry found two accurate and precise in tests that controlled for colour-matching error, four accurate or precise but not both, and two neither~\cite{Reddy2020}.

Applied to that evidence, the framework permits screening under a stated decision rule with confirmation of screen positives and of a sample of negatives by a reference method. At 10~\textmu{}g/L the recorded errors all missed contaminated wells, which is the case for confirming negatives, whereas at 50~\textmu{}g/L they ran in both directions; the threshold and the fraction of negatives confirmed remain the programme's choice, and kit and laboratory costs in Bangladesh were not assessed. The evidence that would change the action is a false-negative rate near the threshold, with its interval, under the rule a programme uses, together with operator and invalid-test rates.

In Ghana the evidence concerns how to organise \textit{E.~coli} testing so that frequency, transport, incubation and staffing remain affordable. Modelled costs per membrane-filtration test for four service approaches, in three study areas in Ghana and Uganda, favoured centralised testing in two of the three~\cite{Trimmer2024}. In the Ghanaian study area, centralised testing was cheapest in 96\% of simulations, at a median of USD 10 per test~\cite{Trimmer2024}. The same study reports the WHO recommendation that samples not be stored for more than 6 h before processing, which links distance to the choice of service, and publishes its cost model for others to run~\cite{Trimmer2024}. A review of commercial field tests found none that met the UNICEF target product profile, with the need for a power source the largest limitation and 18 to 24 h of incubation still required~\cite{Pichel2023}.

The framework treats the instrument as one component of a testing service, and compares service architectures on cost per valid result rather than cost per test. Cost per valid result counts invalid and repeat tests as well, where a valid microbial result is one processed within the holding time, incubated as specified and accompanied by passing controls; this study did not measure it. Until that figure is available, modelled cost per test by service approach, bounded to its study areas and assumptions, is the basis a programme has. Neither costing study states a price year, although one converted costs collected in December 2014 and January 2015 at 1 January 2015 exchange rates~\cite{Delaire2017}, and none of the 23 products in the documentary catalogue documents a per-test cost with its basis, currency and date. Neither setting's evidence supports naming a kit or a product.

\subsection*{Technology evidence and monitoring decisions}

Wherever an action calls for measurement, the choice of instrument rests on the documentation located for 23 commercial products. The 23-product grid holds them against 54 criteria, 1,242 cells (Fig~\ref{fig:product-evidence}). Of these, 891 (71.7\%) are not reported in any located document, 195 (15.7\%) are document-sourced, 9 of them from regulatory certification, 106 were not researched, all for two products, 40 are conflicting, adjacent-product or range-only, and 10 are not applicable. Not reported means that the located documentation is silent: the criterion may still have been measured, and no product has been found wanting. The 9 regulatory cells all belong to one product; they are independent of the manufacturer but are documentation, not device-versus-reference testing. Among the records assembled for the 23 products, paired device-versus-reference field evidence was located only for arsenic, and only one of the three recorded comparisons has recoverable publication provenance.\label{callout:fig:product-evidence}

\begin{figure}[!ht]
\caption{\textbf{Documentary evidence for 23 products across 54 criteria.} Evidence classes of the 1,242 cells: 891 not reported in any located document (unknown evidence), 106 not researched, 40 conflicting, adjacent-product or range-only, 10 not applicable and 195 document-sourced, 9 from regulatory certification. Identities and sources: \nameref{S4_Table}. MTBF, mean time between failures; QC, quality control; CV, coefficient of variation; IP, ingress protection. Data: \nameref{S9_Data}.}
\label{fig:product-evidence}
\end{figure}

Under the two counting rules set out in Methods, 1 of the 11 active criteria meets the documentation threshold on the strict field and 4 on any admissible field (Fig~\ref{fig:documentation-threshold}). None of the 23 products, including each of the 21 that were researched, documents all 11 active criteria. On any admissible field the most any product documents is 9, reached by 1 product, with 4 products documenting 8 and 4 documenting 7; on the strict field the most is 7. Only one criterion therefore supports a comparison across most of the catalogue on its primary field. The outcome therefore rests on requiring all 11: a rule requiring 9 would admit one product. For 19 of the 54 criteria, no located document covers any product. In all 26 contexts the framework named no product, from one determination over the 23-product catalogue that is not evidence about the contexts.\label{callout:fig:documentation-threshold}

\begin{figure}[!ht]
\caption{\textbf{Coverage of the 11 active criteria.} Products documenting each active criterion on strict primary fields and on any admissible field, against the documentation threshold of 12 of 23 products (dashed line). No product documents all 11 active criteria; the most any documents is 9. Data: \nameref{S9_Data}.}
\label{fig:documentation-threshold}
\end{figure}

Among the ten technology archetypes, a population separate from the named products, one archetype led under every cost mapping in 5 of 7 weighting scenarios, and in the other two the leader changed with the cost mapping (\nameref{S7_Table}). A leading archetype is preference context under stated weights, not a product or the class of instrument to use.

Table~\ref{table4} brings the settings together: for each, the evidence held, what a model may and may not do, the action returned and the evidence that would change it; \nameref{S12_Table} adds the monitoring decision each setting poses and where the framework declines.\label{callout:table4}

\nolinenumbers
\begin{spacing}{1}
{\footnotesize\setlength{\tabcolsep}{3pt}\setlength{\LTcapwidth}{\linewidth}
\begin{longtable}{>{\raggedright\arraybackslash}p{0.1323\linewidth}>{\raggedright\arraybackslash}p{0.2523\linewidth}>{\raggedright\arraybackslash}p{0.1800\linewidth}>{\raggedright\arraybackslash}p{0.1800\linewidth}>{\raggedright\arraybackslash}p{0.1923\linewidth}}
\caption{\textbf{Monitoring evidence and actions across application settings.}}\label{table4} \\
\hline
\textbf{Setting} & \textbf{Evidence held} & \textbf{What a model may do} & \textbf{Action returned} & \textbf{Evidence that would change it} \\ \thickhline
\endfirsthead
\multicolumn{5}{@{}l}{\footnotesize\emph{Table 4. (continued)}} \\
\hline
\textbf{Setting} & \textbf{Evidence held} & \textbf{What a model may do} & \textbf{Action returned} & \textbf{Evidence that would change it} \\ \thickhline
\endhead
India (CGWB portal and CGWB Year Book) & 186 cells in 31 states and Union Territories; 160 eligible; 17 decision-relevant at a 0.10 action level, which locate where further records could change a classification and set no action; median revisit interval annual or longer in all 145 cells with one recorded & Ordering may sequence confirmatory sampling for conductivity and TDS (CGWB portal) and fluoride and nitrate (CGWB Year Book); pH not supported; arsenic insufficient support; no pair licenses a probability & Rank prioritisation plus direct confirmation (4); Monitoring baseline required (2); Direct measurement required (1); Targeted direct reference measurement (1) & Local direct measurements permitting recalibration; persistent well, date and sample identity; arsenic method and detection-limit metadata; accepted turbidity and \textit{E.~coli} populations \\ \hline
Kenya, Kwale County (groundwater and springs) & Records above threshold: pH 41 of 140, conductivity 19 of 140, nitrate 5 of 139, fluoride 3 of 139, arsenic 3 of 140; 67 sites sampled in both 2016 campaigns & None; no model was fitted or scored & Local signal plus direct confirmation (5); Monitoring baseline required (3) & Repeat same-site measurement across further seasons of the year; 9 of 24 sites outside the pH range in March were inside it in June \\ \hline
Malawi, laboratory snapshot & Records above threshold: pH 7 of 31, nitrate 0 of 31, turbidity 2 of 22; faecal coliform detectable in 4 of 19; conductivity, fluoride and arsenic not readable as recorded, and TDS on the same mixed scale as conductivity in 10 of 31 records, so no TDS count is given & None; no model was fitted or scored & Local signal plus direct confirmation (4); Data repair required (3); Organism-specific baseline required (2) & Recovered conductivity and TDS scale, fluoride codebook and arsenic units; organism-specific \textit{E.~coli} measurement; wider repeat-site coverage \\ \hline
Malawi, conductivity logger series & 19 sites, 714 site-days over 39 days; 76 site-days above 1,500 \textmu{}S/cm; 2 sites above on every recorded day; 0 changes of state in 695 transitions & None; temporal description only & High-frequency sentinel plus chemical confirmation (1) & Direct chemistry with calibration, fouling, downtime and energy records at the persistent and comparison sites \\ \hline
Bangladesh (published evidence only) & One kit classified 110 of 123 wells correctly against 10 \textmu{}g/L and 113 against the Bangladesh standard of 50 \textmu{}g/L; all 13 errors at 10 \textmu{}g/L were underestimates, and at 50 \textmu{}g/L errors ran in both directions; in tests controlling colour-matching error, two of eight kits were accurate and precise & None; no model transferred & Screening under a stated rule, with reference confirmation of positives and of a sample of negatives; the threshold and the fraction confirmed are the programme's to set & A false-negative rate near the threshold, with its interval, under the programme's rule; operator and invalid-test rates \\ \hline
Ghana (published evidence only) & Modelled cost per test favoured centralised testing in two of three study areas in Ghana and Uganda; in the Ghanaian area it was cheapest in 96\% of simulations; no reviewed field test met the UNICEF target product profile & None; no model transferred & Choose a service architecture on cost per valid result, with the instrument as one component & Network-specific cost per valid result, including invalid and repeat tests, and the sampling frequency at which onsite testing becomes cheaper \\ \hline
\end{longtable}
\begin{flushleft}
1. India, Kenya and Malawi rows summarise the 26 contexts of Table~\ref{table1}; Kwale County and the Malawi archives are local, not national, samples. Bangladesh and Ghana use published evidence only, with no model transferred: one kit at two thresholds, and modelled costs. The full synthesis is \nameref{S12_Table}.

2. Counts above a threshold are record fractions within each archive, not prevalence, and a cell is not a well or a person; thresholds as in Methods (arsenic 0.010~mg/L is 10~\textmu{}g/L), where 1,500~\textmu{}S/cm is an operational project convention, not a health guideline. The Kwale counts exclude surface water and include one thermal spring (Results). Cost per valid result is a requirement, not an outcome measured here. No product is named in any setting.
\end{flushleft}
}
\end{spacing}
\linenumbers

\FloatBarrier

\section*{Discussion}

\subsection*{Implications for monitoring practice}

Taken together, the applications show which kind of evidence carried each action. The 18 Kenyan and Malawian actions came from record harmonisation and data status alone, and in Malawi harmonisation turned unreadable parameters into a defined data-repair task rather than a silent gap. The Indian actions drew on record support and on the transfer verdicts, which permitted ordering in four of the six Indian pairs and supported no probability. The validation designs entered no action; with one partition per design on seven tasks, the fall in mean ROC-AUC from random-row to block splitting illustrates why the framework ties a model's permitted use to the design rather than to the metric. The documentary audit turned missing product information into a criterion-by-criterion account of what no located document reports (\nameref{S4_Table}).

For India, the results permit ordering to sequence confirmation within the four supported archive \texttimes{} parameter pairs, rather than supporting a hazard map. Where ordering transferred, it could decide which conductivity, TDS, nitrate or fluoride sources to confirm first, though not which may be left untested. At the assumed 0.10 action level, the 17 decision-relevant cells locate where further records could change a state-level classification, whereas resolved cells such as Rajasthan call for contaminant-specific confirmation rather than more records. Because calibration slopes were below 1 on every transfer pair and no pair was shown to be adequately calibrated, the framework withholds probability-based use; because the transfer also crossed water-body type, the supported pairs do not separate a change of geography from a change of water type. These implications rest on archived evidence; no operational gain has been tested.

In Kwale County and Malawi the record itself set the next step. In Kwale that step is returning to the same sources, since pH classifications changed between campaigns at 11 of the 67 sites sampled twice, and the fluoride signal rests on a thermal spring and an unverified value that a practitioner would check before acting. In the Malawi laboratory snapshot, four parameters could not be read as recorded, and three of them returned data repair, which needs units, a missing-value codebook and, for arsenic, detection limits to be recovered; these gaps describe the records and say nothing about the water. The snapshot's gaps are in line with reports that groundwater monitoring data in Malawi are persistently scarce, a scarcity made worse by existing monitoring networks that are not operating~\cite{Robertson2026-PLOSW0000593}. The logger series shows persistence at two sites without identifying its cause. Both sites are in Chikwawa District, and groundwater salinity too high for potable use has been documented in Malawi's lower Shire Valley~\cite{BGS-Earthwise-Malawi}; whether that explains these sites was not assessed.

In Bangladesh and Ghana the constraint lay in the decision rule and the service design rather than in any model. In hypothetical well-switching scenarios for Araihazar, Bangladesh, relying on kit rather than laboratory data lost only a modest part of the reduction in average arsenic exposure, set against the logistical and financial challenge of supplying laboratory data alone~\cite{Jameel2021}. The same analysis found the 50~\textmu{}g/L threshold used in Bangladesh, rather than the WHO guideline of 10~\textmu{}g/L, close to optimal for average exposure reduction~\cite{Jameel2021}. A programme could accept kit screening on that basis while setting the residual error it will tolerate. For Ghana, both costing studies cited here, one across sub-Saharan Africa and one covering study areas in Ghana and Uganda, express cost per test with quality-control samples added as a markup~\cite{Delaire2017,Trimmer2024}, whereas a service design needs cost per valid result, counting invalid and repeat tests.

\subsection*{Implications for model use and technology assessment}

Mean discrimination fell from 0.9396 to 0.7696 between two validation designs applied to the same models and data, whereas a review of machine-learning models for United States drinking water found tenfold cross-validation or a randomly chosen test set the most common basis for reported performance~\cite{Hu2023}. In a forest biomass mapping study, non-spatial validation suggested that a model predicted more than half of the variation, whereas spatial validation found almost no predictive power~\cite{Ploton2020}. A methodological source on spatial dependence records the counter-argument that spatial cross-validation is not the right way to evaluate map accuracy~\cite{MeyerPebesma2022}, so the single spatial-block split is used here to address a generalisation claim, not to certify accuracy away from sampled locations. Whether the national hazard models for India cited in the Introduction carry such optimism was not tested here. What the framework contributes here is a rule that ties the validation design to the use a model is allowed, so that a random-row figure cannot authorise decisions in unsampled places.

The calibration result has precedents. Estimated risks can be unreliable even when discrimination is good, for example overestimating risk when a model is used in a setting where incidence is lower~\cite{VanCalster2019}. Calibration-in-the-large, the calibration slope and discrimination have been proposed as distinct measures of model performance, with problems of calibration-in-the-large often found at external validation~\cite{SteyerbergVergouwe2014}, and moderate calibration has been proved to guarantee decision-making that does no harm~\cite{VanCalster2016}. No transferred score here was shown to be moderately calibrated, so none carries that guarantee; the slopes show that transferred probabilities were miscalibrated, not every form the miscalibration takes. The permission gate turns this distinction into a decision rule: a model whose ordering transfers may sequence confirmatory sampling, and only local calibration can license a probability.

Evidence from outside the product catalogue points the same way as the documentary shortfall. Besides the operator-dependent kit verification cited in the Introduction, a laboratory evaluation of a marketed low-cost sensor found no correlation between its score and the presence of \textit{E.~coli}~\cite{Purvis2026-PLOSW0000442}. Both evaluated one product at a stated date, and no conclusion about any product is drawn from them. Most of 72 institutions assessed across ten sub-Saharan African countries did not achieve the microbial testing levels set by applicable standards or WHO guidelines~\cite{Peletz2016}, and drinking-water surveillance often falls short of its required frequency~\cite{GLAAS2025}. Threshold-near performance, operator dependence and achievable testing frequency are what a programme would need documented, by suppliers, service records or independent evaluation, before it could rely on a named instrument; none is among the active criteria audited here.

Abstention here differs from its usual machine-learning form, in which a model abstains where it is likely to make a mistake~\cite{Hendrickx2024}. The framework abstains where the evidence needed for a decision is absent, so its abstention describes the evidence base rather than predictive error. Stochastic multicriteria acceptability analysis aids decisions in problems with inaccurate, uncertain or missing information~\cite{TervonenLahdelma2007}, and the archetype comparison draws on it, but no sampling of weights can supply a criterion value that no located document holds. We therefore treat the documentary audit, not the abstention, as the substantive result: it names the criterion-level documentation the located sources do not report, and the product-level abstention is what the sufficiency rule returns for that evidence base.

\subsection*{Applying the framework in other settings}

Applying the framework elsewhere would need at least the following steps, several of which go beyond what was done here. First, identify and quality-check the local archives and their provenance. Second, align units, analytes and source types, and justify thresholds locally. Third, assess site identity, record support and spatial and temporal dependence before choosing validation designs. Fourth, validate any model's ordering in the new setting and calibrate it on local measurements before its probabilities are used. Fifth, assemble technology and deployment evidence, including threshold-near performance, invalid-test rates and service constraints. Sixth, assess the feasibility of the returned actions with practitioners and test prospectively whether they change monitoring decisions for the better.

Two observations would test the framework itself. It would gain support if the same rules, applied to a catalogue with fuller documentation, admitted a named product, since that would show that the abstention reported here reflects this evidence base rather than rules no product could meet. It would be shown to be underspecified if two analysts applying the stated rules to the same archive returned different actions. No one outside this project has yet applied the rules, the framework has not been compared with an alternative approach, and no monitoring programme has reviewed the returned actions for feasibility or cost.

The evidence that would change the actions reported here is specific: local measurements that allow model recalibration, persistent site and well identity, paired chemistry at sentinel sites, threshold-near kit performance under a stated rule, and invalid and repeat tests recorded in a deployed network, from which cost per valid result follows. \nameref{S8_Table} sets out these five and 11 further research priorities in four families. The measurements a recalibration needs are themselves direct measurements, so its value to a programme depends on the model then reducing later measurement, which was not assessed.

\FloatBarrier

\section*{Conclusion}

The evidence-to-action framework links monitoring records, model evaluation and instrument documentation to monitoring actions, each recorded with the rule and evidence it rests on and the evidence required to reassess it. It was applied, using archived records and published evidence, to 26 monitoring contexts in India, Kwale County (Kenya) and Malawi. Where records alone were available, their state set repeat measurement at the same Kwale County sites and data repair for three Malawian laboratory parameters that could not be read as recorded. In India, 17 of 160 sufficiently sampled state-level parameter cells had intervals straddling the assumed 0.10 action level, so further records could change at least that many classifications. For models, transfer tests supported ordering in four of six Indian archive \texttimes{} parameter pairs, so ordering may sequence confirmatory sampling within those archives; calibration slopes were below 1 on all 13 evaluated pairs, no transferred score was used as a probability, and none was shown to be adequately calibrated. For instruments, no product documented all 11 criteria a comparison requires in any located document, so the audit names, criterion by criterion, what no located document reports.

Applied elsewhere, it would first require verified local records, locally justified thresholds, local tests of any model's ordering and calibration, and deployment evidence for candidate instruments; whether the returned actions improve monitoring remains to be tested prospectively.

\FloatBarrier
\section*{Acknowledgments}

We acknowledge the originators of the water-quality records held in GEMStat and the UNEP GEMS/Water Programme, whose database supplied the records used to fit the models. Indian groundwater-quality data were produced by the Central Ground Water Board and obtained through the National Water Data Portal of the National Water Informatics Centre, Ministry of Jal Shakti, and from the Board's Ground Water Quality Data documents. United States records were obtained from the Water Quality Portal of the National Water Quality Monitoring Council, the U.S. Geological Survey and the U.S. Environmental Protection Agency. The Kenyan and Malawian local archives are openwashdata packages. Soil properties are SoilGrids 2.0 data obtained through ISRIC -- World Soil Information; land cover is ESA WorldCover 10 m 2021 v200; and population density is from WorldPop.

Further acknowledgments: TO BE FILLED LATER

\nolinenumbers

% References: generated with plos2025.bst from references.bib and embedded here, so this file is self-contained.
\begin{thebibliography}{10}

\bibitem{JMP2025}
{WHO/UNICEF Joint Monitoring Programme for Water Supply, Sanitation and
  Hygiene}.
\newblock {Progress on household drinking water, sanitation and hygiene
  2000--2024: special focus on inequalities}.
\newblock Geneva: World Health Organization and United Nations Children's Fund;
  2025.
\newblock ISBN 978-92-4-011514-9.

\bibitem{GLAAS2025}
{World Health Organization}, {United Nations Children's Fund}.
\newblock {State of systems for drinking-water, sanitation and hygiene: global
  update 2025}.
\newblock Geneva: World Health Organization and United Nations Children's Fund;
  2025.
\newblock ISBN 978-92-4-011898-0.

\bibitem{Bain2014}
{Bain} R, {Cronk} R, {Wright} J, {Yang} H, {Slaymaker} T, {Bartram} J.
\newblock {Fecal contamination of drinking-water in low- and middle-income
  countries: a systematic review and meta-analysis}.
\newblock PLoS Med. 2014;11(5):e1001644.
\newblock \href {http://dx.doi.org/10.1371/journal.pmed.1001644}
  {doi:10.1371/journal.pmed.1001644}.

\bibitem{CGWB-SOP-GWQDataAnalysis}
{Central Ground Water Board}.
\newblock {Standard operating procedure on ground water quality data analysis:
  step by step guidelines for ground water quality data analysis}.
\newblock Faridabad: Central Ground Water Board, Ministry of Jal Shakti; [date
  unknown].
\newblock Available from:
  \url{https://cgwb.gov.in/cgwbpnm/public/uploads/documents/17333891661542642909file.pdf}
  [cited 2026 Sep 15].

\bibitem{Podgorski2020}
{Podgorski} J, {Berg} M.
\newblock {Global threat of arsenic in groundwater}.
\newblock Science. 2020;368(6493):845-50.
\newblock \href {http://dx.doi.org/10.1126/science.aba1510}
  {doi:10.1126/science.aba1510}.

\bibitem{Podgorski2018}
{Podgorski} JE, {Labhasetwar} P, {Saha} D, {Berg} M.
\newblock {Prediction modeling and mapping of groundwater fluoride
  contamination throughout India}.
\newblock Environ Sci Technol. 2018;52(17):9889-98.
\newblock \href {http://dx.doi.org/10.1021/acs.est.8b01679}
  {doi:10.1021/acs.est.8b01679}.

\bibitem{Podgorski2020-IJERPH-India}
{Podgorski} J, {Wu} R, {Chakravorty} B, {Polya} DA.
\newblock {Groundwater arsenic distribution in India by machine learning
  geospatial modeling}.
\newblock Int J Environ Res Public Health. 2020;17(19):7119.
\newblock \href {http://dx.doi.org/10.3390/ijerph17197119}
  {doi:10.3390/ijerph17197119}.

\bibitem{Meyer2019}
{Meyer} H, {Reudenbach} C, {W{\"o}llauer} S, {Nauss} T. {Importance of spatial
  predictor variable selection in machine learning applications -- moving from
  data reproduction to spatial prediction} [Preprint]. arXiv:1908.07805v1;
  2019.
\newblock \href {http://dx.doi.org/10.48550/arXiv.1908.07805}
  {doi:10.48550/arXiv.1908.07805}.

\bibitem{MeyerPebesma2022}
{Meyer} H, {Pebesma} E.
\newblock {Machine learning-based global maps of ecological variables and the
  challenge of assessing them}.
\newblock Nat Commun. 2022;13(1):2208.
\newblock \href {http://dx.doi.org/10.1038/s41467-022-29838-9}
  {doi:10.1038/s41467-022-29838-9}.

\bibitem{NiculescuMizil2005}
{Niculescu-Mizil} A, {Caruana} R.
\newblock {Predicting good probabilities with supervised learning}.
\newblock In: Proceedings of the 22nd International Conference on Machine
  Learning (ICML '05). ACM Press; 2005. p. 625-32.
\newblock \href {http://dx.doi.org/10.1145/1102351.1102430}
  {doi:10.1145/1102351.1102430}.

\bibitem{VanCalster2019}
{Van Calster} B, {McLernon} DJ, {van Smeden} M, {Wynants} L, {Steyerberg} EW;
  {Topic Group `Evaluating diagnostic tests and prediction models' of the
  STRATOS initiative}.
\newblock {Calibration: the Achilles heel of predictive analytics}.
\newblock BMC Med. 2019;17(1):230.
\newblock \href {http://dx.doi.org/10.1186/s12916-019-1466-7}
  {doi:10.1186/s12916-019-1466-7}.

\bibitem{Andres2018}
{Andres} L, {Boateng} K, {Borja-Vega} C, {Thomas} E.
\newblock {A review of in-situ and remote sensing technologies to monitor water
  and sanitation interventions}.
\newblock Water (Basel). 2018;10(6):756.
\newblock \href {http://dx.doi.org/10.3390/w10060756} {doi:10.3390/w10060756}.

\bibitem{EPA-ETV-QuickLowRangeII-2003}
{US Environmental Protection Agency Environmental Technology Verification
  Program; Battelle Advanced Monitoring Systems Center}.
\newblock {ETV joint verification statement: arsenic test kit, Quick Low Range
  II (Industrial Test Systems, Inc.)}.
\newblock Washington (DC): US Environmental Protection Agency; 2003 Sep.
\newblock Available from:
  \url{https://archive.epa.gov/nrmrl/archive-etv/web/pdf/01_vs_lowrangell.pdf}
  [cited 2026 Sep 15].

\bibitem{Delaire2017}
{Delaire} C, {Peletz} R, {Kumpel} E, {Kisiangani} J, {Bain} R, {Khush} R.
\newblock {How much will it cost to monitor microbial drinking water quality in
  sub-Saharan Africa?}
\newblock Environ Sci Technol. 2017;51(11):5869-78.
\newblock \href {http://dx.doi.org/10.1021/acs.est.6b06442}
  {doi:10.1021/acs.est.6b06442}.

\bibitem{Schuwirth2018}
{Schuwirth} N, {Honti} M, {Logar} I, {Stamm} C.
\newblock {Multi-criteria decision analysis for integrated water quality
  assessment and management support}.
\newblock Water Res X. 2018;1:100010.
\newblock \href {http://dx.doi.org/10.1016/j.wroa.2018.100010}
  {doi:10.1016/j.wroa.2018.100010}.

\bibitem{Broekhuizen2015}
{Broekhuizen} H, {Groothuis-Oudshoorn} CGM, {van Til} JA, {Hummel} JM,
  {IJzerman} MJ.
\newblock {A review and classification of approaches for dealing with
  uncertainty in multi-criteria decision analysis for healthcare decisions}.
\newblock Pharmacoeconomics. 2015;33(5):445-55.
\newblock \href {http://dx.doi.org/10.1007/s40273-014-0251-x}
  {doi:10.1007/s40273-014-0251-x}.

\bibitem{Hendrickx2024}
{Hendrickx} K, {Perini} L, {Van der Plas} D, {Meert} W, {Davis} J.
\newblock {Machine learning with a reject option: a survey}.
\newblock Mach Learn. 2024;113(5):3073-110.
\newblock \href {http://dx.doi.org/10.1007/s10994-024-06534-x}
  {doi:10.1007/s10994-024-06534-x}.

\bibitem{Hu2023}
{Hu} XC, {Dai} M, {Sun} JM, {Sunderland} EM.
\newblock {The utility of machine learning models for predicting chemical
  contaminants in drinking water: promise, challenges, and opportunities}.
\newblock Curr Environ Health Rep. 2023;10(1):45-60.
\newblock \href {http://dx.doi.org/10.1007/s40572-022-00389-x}
  {doi:10.1007/s40572-022-00389-x}.

\bibitem{Geifman2017}
{Geifman} Y, {El-Yaniv} R.
\newblock {Selective classification for deep neural networks}.
\newblock In: {Guyon} I, {Von Luxburg} U, {Bengio} S, {Wallach} H, {Fergus} R,
  {Vishwanathan} S, et~al., editors. Advances in Neural Information Processing
  Systems 30 (NIPS 2017). Curran Associates; 2017. Available from:
  \url{https://proceedings.neurips.cc/paper_files/paper/2017/file/4a8423d5e91fda00bb7e46540e2b0cf1-Paper.pdf}
  [cited 2026 Sep 16].

\bibitem{GEMStat}
{UNEP GEMS/Water Programme}.
\newblock {GEMStat database} [Internet]; [date unknown] [cited 2026 Sep 1].
\newblock Available from: \url{https://gemstat.org/data-gemstat/}.

\bibitem{CGWB-NWDP}
{Central Ground Water Board}.
\newblock {Ground Water Quality (Manual - Chemical Parameters), Central Ground
  Water Board (CGWB)} [dataset]. National Water Data Portal, National Water
  Informatics Centre, Ministry of Jal Shakti; [date unknown] [cited 2026 Aug
  7].
\newblock Available from:
  \url{https://nwdp.nwic.gov.in/dataset/7bb1e7c7-bcc1-48bb-8bcd-32c19473804a}.

\bibitem{CGWB-GWQ2019}
{Central Ground Water Board}.
\newblock {Ground Water Quality Data -- 2019} [data tables, PDF]. Ministry of
  Jal Shakti; [date unknown] [cited 2026 Sep 12].
\newblock Available from:
  \url{https://cgwb.gov.in/cgwbpnm/public/uploads/documents/16860545111615830419file.pdf}.
\newblock Internet Archive snapshot 20240124163108.

\bibitem{CGWB-GWQ2021}
{Central Ground Water Board}.
\newblock {Ground Water Quality Data -- 2021} [data tables, PDF]. Ministry of
  Jal Shakti; [date unknown] [cited 2026 Sep 12].
\newblock Available from:
  \url{https://cgwb.gov.in/cgwbpnm/public/uploads/documents/1686055122486176238file.pdf}.
\newblock Internet Archive snapshot 20231015091851.

\bibitem{CGWB-GWQ2022}
{Central Ground Water Board}.
\newblock {Ground Water Quality Data -- 2022 (Ground Water Quality of
  Unconfined Aquifers)} [data tables, PDF]. Ministry of Jal Shakti; [date
  unknown] [cited 2026 Sep 12].
\newblock Available from:
  \url{https://cgwb.gov.in/sites/default/files/inline-files/final_nhs-wq_pre_2022_compressed.pdf}.

\bibitem{CGWB-AGWQR2025}
{Central Ground Water Board}.
\newblock {Annual ground water quality report-2025}.
\newblock Faridabad: Central Ground Water Board, Ministry of Jal Shakti,
  Government of India; 2025 Nov.
\newblock Available from:
  \url{https://cgwb.gov.in/cgwbpnm/public/uploads/documents/1762854375262680475file.pdf}
  [cited 2026 Sep 16].

\bibitem{WQP}
{Water Quality Portal}. Washington (DC): National Water Quality Monitoring
  Council, United States Geological Survey (USGS), Environmental Protection
  Agency (EPA); 2021.
\newblock \href {http://dx.doi.org/10.5066/P9QRKUVJ} {doi:10.5066/P9QRKUVJ}.

\bibitem{OWD-grdwtrsmpkwale}
{Loos} S, {Zhong} M, {Hope} R. {grdwtrsmpkwale: Groundwater analysis from 2016
  in Kwale, Kenya} [dataset]. Zenodo; 2023.
\newblock Version 0.0.1.
\newblock \href {http://dx.doi.org/10.5281/zenodo.10016875}
  {doi:10.5281/zenodo.10016875}.

\bibitem{OWD-boreholelabdata}
{Mhango} E, {Paterson} F, {Rattray} J. {boreholelabdata: Water Quality
  Laboratory Analysis and Reporting Dataset -- Malawi (2017--2019)} [dataset].
  Zenodo; 2025.
\newblock \href {http://dx.doi.org/10.5281/zenodo.17433461}
  {doi:10.5281/zenodo.17433461}.

\bibitem{OWD-mwgroundwaterdata}
{Mhango} E. {mwgroundwaterdata: Malawi Groundwater Monitoring Time-Series
  Dataset (2024--2025)} [dataset]. Zenodo; 2026.
\newblock \href {http://dx.doi.org/10.5281/zenodo.18798409}
  {doi:10.5281/zenodo.18798409}.

\bibitem{WHO2022GDWQ}
{World Health Organization}.
\newblock {Guidelines for drinking-water quality: fourth edition incorporating
  the first and second addenda}.
\newblock Geneva: World Health Organization; 2022.
\newblock ISBN 978-92-4-004506-4.

\bibitem{USEPA-NPDWR}
{United States Environmental Protection Agency}.
\newblock {National Primary Drinking Water Regulations} [Internet]. Washington
  (DC): US EPA; [date unknown] [updated 2026 Aug 31; cited 2026 Sep 15].
\newblock Available from:
  \url{https://www.epa.gov/ground-water-and-drinking-water/national-primary-drinking-water-regulations}.

\bibitem{Brown2001}
{Brown} LD, {Cai} TT, {DasGupta} A.
\newblock {Interval estimation for a binomial proportion}.
\newblock Stat Sci. 2001;16(2):101-33.
\newblock \href {http://dx.doi.org/10.1214/ss/1009213286}
  {doi:10.1214/ss/1009213286}.

\bibitem{Pedregosa2011}
{Pedregosa} F, {Varoquaux} G, {Gramfort} A, {Michel} V, {Thirion} B, {Grisel}
  O, et~al.
\newblock {Scikit-learn: machine learning in Python}.
\newblock J Mach Learn Res. 2011;12:2825-30.
\newblock Available from:
  \url{https://www.jmlr.org/papers/v12/pedregosa11a.html}.

\bibitem{Saito2015}
{Saito} T, {Rehmsmeier} M.
\newblock {The precision-recall plot is more informative than the ROC plot when
  evaluating binary classifiers on imbalanced datasets}.
\newblock PLoS One. 2015;10(3):e0118432.
\newblock \href {http://dx.doi.org/10.1371/journal.pone.0118432}
  {doi:10.1371/journal.pone.0118432}.

\bibitem{CGWB-ManualAquiferMapping}
{Central Ground Water Board}.
\newblock {Manual on aquifer mapping}.
\newblock Faridabad: Central Ground Water Board, Ministry of Water Resources;
  [date unknown].
\newblock Available from:
  \url{https://cgwb.gov.in/cgwbpnm/public/uploads/documents/16861227981023889265file.pdf}
  [cited 2026 Sep 15].

\bibitem{Poggio2021}
{Poggio} L, {de Sousa} LM, {Batjes} NH, {Heuvelink} GBM, {Kempen} B, {Ribeiro}
  E, et~al.
\newblock {SoilGrids 2.0: producing soil information for the globe with
  quantified spatial uncertainty}.
\newblock SOIL. 2021;7(1):217-40.
\newblock \href {http://dx.doi.org/10.5194/soil-7-217-2021}
  {doi:10.5194/soil-7-217-2021}.

\bibitem{Zanaga2022}
{Zanaga} D, {Van De Kerchove} R, {Daems} D, {De Keersmaecker} W, {Brockmann} C,
  {Kirches} G, et~al.
\newblock {ESA WorldCover 10 m 2021 v200 [dataset]}.
\newblock Zenodo; 2022.
\newblock \href {http://dx.doi.org/10.5281/zenodo.7254221}
  {doi:10.5281/zenodo.7254221}.

\bibitem{WorldPop2020}
{WorldPop}, {Bondarenko} M. {Individual countries 1km population density
  (2000--2020)} [dataset]. University of Southampton; 2020.
\newblock \href {http://dx.doi.org/10.5258/SOTON/WP00674}
  {doi:10.5258/SOTON/WP00674}.

\bibitem{Chakraborty2022}
{Chakraborty} S.
\newblock {TOPSIS and modified TOPSIS: a comparative analysis}.
\newblock Decis Anal J. 2022;2:100021.
\newblock \href {http://dx.doi.org/10.1016/j.dajour.2021.100021}
  {doi:10.1016/j.dajour.2021.100021}.

\bibitem{TervonenLahdelma2007}
{Tervonen} T, {Lahdelma} R.
\newblock {Implementing stochastic multicriteria acceptability analysis}.
\newblock Eur J Oper Res. 2007;178(2):500-13.
\newblock \href {http://dx.doi.org/10.1016/j.ejor.2005.12.037}
  {doi:10.1016/j.ejor.2005.12.037}.

\bibitem{BGS-Earthwise-Malawi}
{Upton} K, {{\'O} Dochartaigh} B{\'E}, {Chunga} B, {Bellwood-Howard} I.
\newblock {Africa Groundwater Atlas: hydrogeology of Malawi} [Internet].
  British Geological Survey; 2018 [updated 2023 Dec 20; cited 2026 Sep 15].
\newblock Available from:
  \url{https://earthwise.bgs.ac.uk/index.php/Hydrogeology_of_Malawi}.

\bibitem{George2012}
{George} CM, {Zheng} Y, {Graziano} JH, {Rasul} SB, {Hossain} Z, {Mey} JL,
  et~al.
\newblock {Evaluation of an arsenic test kit for rapid well screening in
  Bangladesh}.
\newblock Environ Sci Technol. 2012;46(20):11213-9.
\newblock \href {http://dx.doi.org/10.1021/es300253p} {doi:10.1021/es300253p}.

\bibitem{Reddy2020}
{Reddy} RR, {Rodriguez} GD, {Webster} TM, {Abedin} MJ, {Karim} MR, {Raskin} L,
  et~al.
\newblock {Evaluation of arsenic field test kits for drinking water:
  recommendations for improvement and implications for arsenic affected regions
  such as Bangladesh}.
\newblock Water Res. 2020;170:115325.
\newblock \href {http://dx.doi.org/10.1016/j.watres.2019.115325}
  {doi:10.1016/j.watres.2019.115325}.

\bibitem{Trimmer2024}
{Trimmer} JT, {Delaire} C, {Marshall} K, {Khush} R, {Peletz} R.
\newblock {Centralized or onsite testing? Examining the costs of water quality
  monitoring in rural Africa}.
\newblock Environ Sci Technol. 2024;58(26):11236-46.
\newblock \href {http://dx.doi.org/10.1021/acs.est.4c01916}
  {doi:10.1021/acs.est.4c01916}.

\bibitem{Pichel2023}
{Pichel} N, {Hymn{\^o} de Souza} F, {Sabogal-Paz} LP, {Shah} PK, {Adhikari} N,
  {Pandey} S, et~al.
\newblock {Field-testing solutions for drinking water quality monitoring in
  low- and middle-income regions and case studies from Latin American, African
  and Asian countries}.
\newblock J Environ Chem Eng. 2023;11(6):111180.
\newblock \href {http://dx.doi.org/10.1016/j.jece.2023.111180}
  {doi:10.1016/j.jece.2023.111180}.

\bibitem{Robertson2026-PLOSW0000593}
{Robertson} DJC, {Byrns} S, {Ellis} R, {White} CJ, {Nhlema} M, {Ngwira} G,
  et~al.
\newblock {Persistent groundwater data scarcity: challenges and advancements
  from a co-creation process in Malawi}.
\newblock PLOS Water. 2026;5(9):e0000593.
\newblock \href {http://dx.doi.org/10.1371/journal.pwat.0000593}
  {doi:10.1371/journal.pwat.0000593}.

\bibitem{Jameel2021}
{Jameel} Y, {Mozumder} MRH, {van Geen} A, {Harvey} CF.
\newblock {Well-switching to reduce arsenic exposure in Bangladesh: making the
  most of inaccurate field kit measurements}.
\newblock Geohealth. 2021;5(12):e2021GH000464.
\newblock \href {http://dx.doi.org/10.1029/2021GH000464}
  {doi:10.1029/2021GH000464}.

\bibitem{Ploton2020}
{Ploton} P, {Mortier} F, {R{\'e}jou-M{\'e}chain} M, {Barbier} N, {Picard} N,
  {Rossi} V, et~al.
\newblock {Spatial validation reveals poor predictive performance of
  large-scale ecological mapping models}.
\newblock Nat Commun. 2020;11(1):4540.
\newblock \href {http://dx.doi.org/10.1038/s41467-020-18321-y}
  {doi:10.1038/s41467-020-18321-y}.

\bibitem{SteyerbergVergouwe2014}
{Steyerberg} EW, {Vergouwe} Y.
\newblock {Towards better clinical prediction models: seven steps for
  development and an ABCD for validation}.
\newblock Eur Heart J. 2014;35(29):1925-31.
\newblock \href {http://dx.doi.org/10.1093/eurheartj/ehu207}
  {doi:10.1093/eurheartj/ehu207}.

\bibitem{VanCalster2016}
{Van Calster} B, {Nieboer} D, {Vergouwe} Y, {De Cock} B, {Pencina} MJ,
  {Steyerberg} EW.
\newblock {A calibration hierarchy for risk models was defined: from utopia to
  empirical data}.
\newblock J Clin Epidemiol. 2016;74:167-76.
\newblock \href {http://dx.doi.org/10.1016/j.jclinepi.2015.12.005}
  {doi:10.1016/j.jclinepi.2015.12.005}.

\bibitem{Purvis2026-PLOSW0000442}
{Purvis} T, {Nguyen} T, {McHugh} C, {Yeolekar} S, {Ding} E, {Wang} J, et~al.
\newblock {The Lishtot TestDrop Pro device cannot reliably indicate microbial
  water safety}.
\newblock PLOS Water. 2026;5(6):e0000442.
\newblock \href {http://dx.doi.org/10.1371/journal.pwat.0000442}
  {doi:10.1371/journal.pwat.0000442}.

\bibitem{Peletz2016}
{Peletz} R, {Kumpel} E, {Bonham} M, {Rahman} Z, {Khush} R.
\newblock {To what extent is drinking water tested in sub-Saharan Africa? A
  comparative analysis of regulated water quality monitoring}.
\newblock Int J Environ Res Public Health. 2016;13(3):275.
\newblock \href {http://dx.doi.org/10.3390/ijerph13030275}
  {doi:10.3390/ijerph13030275}.

\end{thebibliography}

\linenumbers
\section*{Supporting information}

\paragraph*{S1 Text.}
\phantomsection
\label{S1_Text}
\textbf{Harmonisation, model configuration and decision rules.} Data rules, configurations, decision-rule branches and the scope of result verification.

\paragraph*{S2 Table.}
\phantomsection
\label{S2_Table}
\textbf{The 26 monitoring contexts in full.} The complete context matrix and the eight action definitions.

\paragraph*{S3 Table.}
\phantomsection
\label{S3_Table}
\textbf{Support, cadence and exceedance for 186 Indian cells.} Per-cell values behind Table 2, Fig~\ref{fig:india-map}, Fig~\ref{fig:india-intervals} and Fig~\ref{fig:india-support-cadence}.

\paragraph*{S4 Table.}
\phantomsection
\label{S4_Table}
\textbf{The 23 \texttimes{} 54 product documentation grid.} Evidence class, value, source and locus for each cell; products P01--P23 as in Fig~\ref{fig:product-evidence}.

\paragraph*{S5 Text.}
\phantomsection
\label{S5_Text}
\textbf{Mechanistic ablation.}

\paragraph*{S6 Fig.}
\phantomsection
\label{S6_Fig}
\textbf{An illustrative conductivity prediction surface.} Read as an ordering, not a probability.

\paragraph*{S7 Table.}
\phantomsection
\label{S7_Table}
\textbf{Rank acceptability of technology archetypes in seven weighting scenarios.} Archetypes are not products.

\paragraph*{S8 Table.}
\phantomsection
\label{S8_Table}
\textbf{Sixteen research priorities in four families.}

\paragraph*{S9 Data.}
\phantomsection
\label{S9_Data}
\textbf{The values plotted in the main-text figures and S6 Fig.} One CSV file per figure.

\paragraph*{S10 Text.}
\phantomsection
\label{S10_Text}
\textbf{Record of artificial-intelligence assistance.}

\paragraph*{S11 Table.}
\phantomsection
\label{S11_Table}
\textbf{Record fractions and campaign comparisons for the Kenyan and Malawian archives.} The archives are local, not national samples.

\paragraph*{S12 Table.}
\phantomsection
\label{S12_Table}
\textbf{The setting synthesis in full.} The full synthesis has seven columns; Table 4 shows five.

\end{document}
